\chapter{Spectral Algebraic Geometry from Commutative
\texorpdfstring{$\infty$}{infinity}-Operadic Algebra}
\label{chap:spectral-algebraic-geometry}
\label{chap:connective-spectral-algebraic-geometry}

\section{Purpose and standing conventions}
\label{sec:spectral-ag-purpose}

The purpose of this chapter is to construct spectral algebraic geometry
from the commutative $\infty$-operad, its interpretation in spectra, and
the canonical Postnikov $t$-structure. Throughout this work, coordinate
algebras are connective commutative ring spectra. Accordingly, an affine
spectral scheme means an object of
\[
    \Aff_{\Sp}
    :=
    \bigl(\CAlg(\Sp)^{\mathrm{cn}}\bigr)^{\op},
\]
and a spectral scheme is locally represented by such affine spectral
schemes. The adjective ``connective'' is therefore suppressed from the
geometric terminology, though it remains explicit when contrasting
coordinate algebras with arbitrary objects of $\CAlg(\Sp)$.

The stable semantic category is
\[
    \Sp^\otimes
    \in
    \CAlg(\Pr^L_{\mathrm{st}}),
\]
and commutative ring spectra are interpreted operadically by
\[
    \CAlg(\Sp)
    =
    \Alg_{\Comm^\otimes}(\Sp^\otimes).
\]
Connectivity is imposed on coordinate algebras and selected modules, not
on the ambient stable categories. In particular, for
$A\in\CAlg(\Sp)^{\mathrm{cn}}$ one retains the full stable module category
\[
    \QCoh(\Spec(A))
    :=
    \Mod_A(\Sp).
\]

The operadic module construction gives a functor
\[
    \mathbf{Mod}_{\Sp}
    \colon
    \CAlg(\Sp)
    \longrightarrow
    \CAlg(\Pr^L_{\mathrm{st}}),
    \qquad
    A
    \longmapsto
    \Mod_A(\Sp),
\]
whose transition functors are extension of scalars. The Postnikov
$t$-structures on spectra and module spectra distinguish geometric
flatness from formal stable exactness: extension of scalars is always
exact between stable $\infty$-categories, while it is flat precisely when
it is $t$-exact.

This chapter constructs affine and global spectral schemes, the spectral
Zariski topology, quasi-coherent modules and algebras, relative spectra,
open and closed geometry, closed gluing, vector bundles, relative linear
spaces, and the brave new affine spaces. Cotangent complexes and
infinitesimal deformation theory are developed in
\Ref{chap:spectral-deformation-theory}; smooth, \`etale, and Nisnevich
geometry are developed in \Ref{chap:smooth-spectral-motivic-site}.

The basic variance is recorded by the definitions and functors
\[
    \CAlg(\Sp)^{\mathrm{cn}}
    =
    \Alg_{\Comm^\otimes}(\Sp^\otimes)^{\mathrm{cn}},
    \qquad
    \Aff_{\Sp}
    :=
    \bigl(\CAlg(\Sp)^{\mathrm{cn}}\bigr)^{\op},
\]
\[
    h
    \colon
    \Aff_{\Sp}
    \hookrightarrow
    \PSh(\Aff_{\Sp}),
\]
and, for every spectral scheme $X$,
\[
    \Spec_X
    \colon
    \operatorname{QCAlg}(X)^{\mathrm{cn},\op}
    \longrightarrow
    \operatorname{Sch}^{\mathrm{sp}}_{/X}.
\]
There is no covariant functor which merely ``takes opposites'' from
$\CAlg(\Sp)^{\mathrm{cn}}$ to $\Aff_{\Sp}$; the opposite category is part
of the definition of the affine geometric category.

The foundational sources are
\cite{LurieHigherAlgebra,LurieSAG,LurieDAGIV}. The conventions below are
used throughout the three foundational chapters.

\begin{convention}
\label{conv:spectral-infinity-categorical}
All categories are $\infty$-categories unless explicitly stated otherwise.
Mapping spaces are denoted by
\[
    \Map_{\mathcal C}(X,Y).
\]
Mapping spectra in a stable $\infty$-category are denoted by
\[
    \map_{\mathcal C}(X,Y).
\]
Equivalences are denoted by
\[
    \simeq.
\]
\end{convention}

\begin{convention}[Higher coherence]
\label{conv:spectral-higher-coherence}
Associativity, unitality, commutativity, operadic, base-change, descent,
and functoriality data are understood together with all required coherent
higher homotopies. The coherent data used below are supplied by the
operadic algebra and module constructions, coCartesian module fibrations,
universal properties, Kan extensions, and descent totalizations. Once
supplied, they are suppressed from the notation.
\end{convention}

\begin{convention}
\label{conv:size-spectral-ag}
All constructions are made relative to fixed Grothendieck universes
\[
    \mathbb U
    \in
    \mathbb V.
\]
Categories of algebras, modules, spectral schemes, and smooth spectral
schemes are formed in $\mathbb U$ whenever possible, while presheaf,
sheaf, and functor categories on essentially $\mathbb U$-small sources
are formed in the larger universe $\mathbb V$. Universe changes are
otherwise suppressed from the notation.

\end{convention}

\begin{convention}[Sheaves and hyperdescent]
\label{conv:spectral-sheaves-hyperdescent}
Unless explicitly stated otherwise,
\[
    \Sh_\tau(\mathcal C)
\]
denotes the $\infty$-category of $\tau$-sheaves of spaces without
hypercompletion. When hyperdescent is required, the hypercomplete
$\infty$-topos is denoted by
\[
    \Sh_\tau(\mathcal C)^\wedge.
\]
\end{convention}

\begin{convention}[Monoidal structures on categories]
\label{conv:spectral-monoidal-structures-on-categories}
The notation
\[
    \CAlg(\Cat_\infty)
\]
uses the Cartesian symmetric monoidal structure on $\Cat_\infty$.
The notation
\[
    \CAlg(\Pr^L_{\mathrm{st}})
\]
uses the tensor product of presentable stable $\infty$-categories. Its
morphisms are colimit-preserving strong symmetric monoidal functors.
\end{convention}

\begin{convention}[Ordinary schemes as spectral schemes]
\label{conv:ordinary-schemes-as-spectral-schemes}
The Eilenberg--Mac Lane functor induces a fully faithful functor
\[
    \operatorname{AffSch}
    \hookrightarrow
    \Aff_{\Sp},
    \qquad
    \Spec(R)
    \longmapsto
    \Spec(HR).
\]
By Zariski gluing, it induces a fully faithful functor
\[
    \operatorname{Sch}
    \hookrightarrow
    \operatorname{Sch}^{\mathrm{sp}}.
\]
Whenever an ordinary scheme occurs as a spectral scheme or spectral
prestack, it is understood through this fully faithful discrete
enhancement. In particular, for an ordinary commutative ring \(R\), the
notation
\[
    \Spec(R)
\]
may denote the ordinary affine scheme or its associated discrete affine
spectral scheme
\[
    \Spec(HR),
\]
according to the ambient category.
\end{convention}

\section{Spectra and the canonical Postnikov structure}
\label{sec:spectra-postnikov-structure}

\subsection{The stable symmetric monoidal category of spectra}

\begin{proposition}[Spectral semantic category]
\label{prop:spectra-presentably-stable-monoidal}
The $\infty$-category of spectra admits a canonical structure
\[
    \Sp^\otimes
    \in
    \CAlg(\Pr^L_{\mathrm{st}}).
\]
Its tensor product is the smash product
\[
    -\otimes-
    \colon
    \Sp\times\Sp
    \longrightarrow
    \Sp,
\]
its tensor unit is the sphere spectrum $\mathbb S$, and its tensor product
preserves small colimits separately in each variable.
\end{proposition}

\begin{proof}
The stabilization of pointed spaces is a presentable stable
$\infty$-category. The smash product on pointed spaces stabilizes to a
symmetric monoidal product on spectra, and the stabilization universal
property identifies its tensor unit with the sphere spectrum. Separate
preservation of colimits follows from the presentably symmetric monoidal
construction; see \cite{LurieHigherAlgebra}.
\end{proof}

\begin{proposition}[Closedness of spectra]
\label{prop:spectra-closedness}
For every spectrum $X$, the functor
\[
    -\otimes X
    \colon
    \Sp
    \longrightarrow
    \Sp
\]
admits a right adjoint
\[
    \HHom_{\Sp}(X,-)
    \colon
    \Sp
    \longrightarrow
    \Sp.
\]
Thus there are natural equivalences
\[
    \Map_{\Sp}(Z\otimes X,Y)
    \simeq
    \Map_{\Sp}
    \bigl(
        Z,
        \HHom_{\Sp}(X,Y)
    \bigr).
\]
\end{proposition}

\begin{proof}
The functor $-\otimes X$ preserves small colimits by
\Ref{prop:spectra-presentably-stable-monoidal}. Since $\Sp$ is
presentable, the adjoint functor theorem supplies the right adjoint.
\end{proof}

The adjunction is displayed by
\[
\begin{tikzcd}[column sep=large]
    \Sp
        \arrow[r, shift left=1.1ex, "-\otimes X"]
    &
    \Sp
        \arrow[l, shift left=1.1ex, "{\HHom_{\Sp}(X,-)}"'].
\end{tikzcd}
\]

\subsection{The canonical \texorpdfstring{$t$}{t}-structure}

\begin{definition}[Postnikov $t$-structure]
\label{def:postnikov-t-structure}
For every integer $n$, let
\[
    \Sp_{\geq n}
\]
be the full subcategory of spectra $X$ for which
\[
    \pi_i(X)=0
    \qquad
    (i<n),
\]
and let
\[
    \Sp_{\leq n}
\]
be the full subcategory of spectra $X$ for which
\[
    \pi_i(X)=0
    \qquad
    (i>n).
\]
The pair
\[
    (\Sp_{\geq0},\Sp_{\leq0})
\]
is the \textbf{canonical Postnikov $t$-structure} on $\Sp$. Its truncation
functors are denoted
\[
    \tau_{\geq n}
    \colon
    \Sp
    \longrightarrow
    \Sp_{\geq n}
\]
and
\[
    \tau_{\leq n}
    \colon
    \Sp
    \longrightarrow
    \Sp_{\leq n}.
\]
\end{definition}

\begin{proposition}[Postnikov truncation triangle]
\label{prop:postnikov-truncation-triangle}
For every spectrum $X$ and every integer $n$, there is a natural fibre and
cofibre sequence
\[
    \tau_{\geq n+1}X
    \longrightarrow
    X
    \longrightarrow
    \tau_{\leq n}X.
\]
The functors $\tau_{\geq n}$ and $\tau_{\leq n}$ are respectively right
and left adjoint to the inclusions of the corresponding full
subcategories.
\end{proposition}

\begin{proof}
This is the defining truncation sequence of the Postnikov $t$-structure.
The equality of fibre and cofibre sequences follows from stability.
\end{proof}

\begin{proposition}[Completeness and detection]
\label{prop:postnikov-completeness-detection}
The Postnikov $t$-structure is accessible, compatible with filtered
colimits, left complete, and right complete. In particular, there are
natural equivalences
\[
    X
    \simeq
    \varprojlim_n\tau_{\leq n}X
\]
and
\[
    X
    \simeq
    \varinjlim_n\tau_{\geq -n}X.
\]
A morphism of spectra is an equivalence if and only if it induces an
isomorphism on every homotopy group.
\end{proposition}

\begin{proof}
These are the standard completeness and accessibility properties of the
Postnikov $t$-structure; see \cite{LurieHigherAlgebra}.
\end{proof}

\begin{definition}[Discrete spectrum]
\label{def:discrete-spectrum}
A spectrum $X$ is \textbf{discrete} if
\[
    X\in\Sp_{\geq0}\cap\Sp_{\leq0}.
\]
The heart of the Postnikov $t$-structure is denoted
\[
    \Sp^\heartsuit.
\]
\end{definition}

\begin{proposition}[Eilenberg--Mac Lane equivalence]
\label{prop:eilenberg-mac-lane-heart}
The Eilenberg--Mac Lane functor induces an equivalence
\[
    H
    \colon
    \operatorname{Ab}
    \xrightarrow{\simeq}
    \Sp^\heartsuit.
\]
Under this equivalence, the inverse functor is $\pi_0$.
\end{proposition}

\begin{proof}
A discrete spectrum is determined by its degree-zero homotopy group, and
every abelian group is realized by an Eilenberg--Mac Lane spectrum.
\end{proof}

\subsection{Compatibility with the smash product}

\begin{proposition}[Connectivity of the smash product]
\label{prop:smash-product-connectivity}
If
\[
    X\in\Sp_{\geq m}
    \qquad\text{and}\qquad
    Y\in\Sp_{\geq n},
\]
then
\[
    X\otimes Y
    \in
    \Sp_{\geq m+n}.
\]
In particular, $\Sp_{\geq0}$ is closed under the smash product and contains
the tensor unit $\mathbb S$.
\end{proposition}

\begin{proof}
After suspension, it is enough to consider $m=n=0$. Connective spectra are
generated under colimits and extensions by suspension spectra of pointed
spaces in nonnegative degrees. The smash product preserves colimits
separately and carries the corresponding generators to connective
spectra. Equivalently, the compatibility follows from the monoidal
compatibility of the Postnikov $t$-structure; see
\cite{LurieHigherAlgebra}.
\end{proof}

\begin{corollary}[Right $t$-exactness of connective tensoring]
\label{cor:connective-tensor-right-t-exact}
For every connective spectrum $X$, the functor
\[
    -\otimes X
    \colon
    \Sp
    \longrightarrow
    \Sp
\]
is right $t$-exact.
\end{corollary}

\begin{proof}
If $Y$ is connective, then $Y\otimes X$ is connective by
\Ref{prop:smash-product-connectivity}.
\end{proof}

\section{Commutative \texorpdfstring{$\infty$}{infinity}-operadic algebra
in spectra}
\label{sec:commutative-operadic-spectral-algebra}

Let
\[
    \Comm^\otimes
    \longrightarrow
    N(\mathrm{Fin}_*)
\]
denote the commutative $\infty$-operad.

\begin{definition}[Commutative ring spectrum]
\label{def:commutative-ring-spectrum}
A \textbf{commutative ring spectrum} is an algebra over the commutative
$\infty$-operad in spectra:
\[
    A
    \in
    \Alg_{\Comm^\otimes}(\Sp^\otimes).
\]
The resulting $\infty$-category is denoted
\[
    \CAlg(\Sp)
    :=
    \Alg_{\Comm^\otimes}(\Sp^\otimes).
\]
Its objects are also called $\EE_\infty$-rings.
\end{definition}

\begin{remark}[Symmetric monoidal envelope]
\label{rem:spectral-commutative-envelope}
Let
\[
    \operatorname{Env}(\Comm)^\otimes
\]
denote the symmetric monoidal envelope of the commutative
$\infty$-operad. Its universal property gives a natural equivalence
\[
    \CAlg(\Sp)
    \simeq
    \Fun^\otimes
    \bigl(
        \operatorname{Env}(\Comm)^\otimes,
        \Sp^\otimes
    \bigr).
\]
Thus the commutative algebraic operations are retained together with all
higher coherence data. See
\cite[Proposition~2.2.4.9]{LurieHigherAlgebra}.
\end{remark}

The operadic and symmetric monoidal-envelope descriptions are related by
\[
\begin{tikzcd}[column sep=huge]
    \Alg_{\Comm^\otimes}(\Sp^\otimes)
        \arrow[r, "\simeq"]
    &
    \Fun^\otimes
    \bigl(
        \operatorname{Env}(\Comm)^\otimes,
        \Sp^\otimes
    \bigr).
\end{tikzcd}
\]

\begin{proposition}[Presentability of commutative ring spectra]
\label{prop:calg-spectra-presentable}
The $\infty$-category
\[
    \CAlg(\Sp)
\]
is presentable. Its forgetful functor
\[
    U
    \colon
    \CAlg(\Sp)
    \longrightarrow
    \Sp
\]
is accessible, creates limits and sifted colimits, and admits a left
adjoint.
\end{proposition}

\begin{proof}
This is the presentability theorem for algebras over an accessible
$\infty$-operad in a presentably symmetric monoidal $\infty$-category
whose tensor product preserves small colimits separately; see
\cite{LurieHigherAlgebra}.
\end{proof}

\subsection{Free commutative ring spectra}

\begin{definition}[Free commutative algebra]
\label{def:free-commutative-spectral-algebra}
The left adjoint to the forgetful functor is denoted
\[
    \operatorname{Sym}
    \colon
    \Sp
    \longrightarrow
    \CAlg(\Sp).
\]
It is the \textbf{free commutative algebra functor} associated to the
commutative $\infty$-operad.
\end{definition}

The adjunction is displayed by
\[
\begin{tikzcd}[column sep=large]
    \Sp
        \arrow[r, shift left=1.1ex, "\operatorname{Sym}"]
    &
    \CAlg(\Sp)
        \arrow[l, shift left=1.1ex, "U"'].
\end{tikzcd}
\]

\begin{proposition}[Underlying free algebra]
\label{prop:underlying-free-spectral-algebra}
For every spectrum $M$, there is a natural equivalence of underlying
spectra
\[
    U\operatorname{Sym}(M)
    \simeq
    \bigoplus_{n\geq0}
    \left(M^{\otimes n}\right)_{h\Sigma_n},
\]
where the summand for $n=0$ is the sphere spectrum.
\end{proposition}

\begin{proof}
This is the explicit free-algebra formula for the commutative
$\infty$-operad in a presentably symmetric monoidal $\infty$-category; see
\cite{LurieHigherAlgebra}.
\end{proof}

\begin{corollary}[Connectivity of free algebras]
\label{cor:free-algebra-connective}
If $M$ is connective, then $\operatorname{Sym}(M)$ is connective.
\end{corollary}

\begin{proof}
Every tensor power $M^{\otimes n}$ is connective by
\Ref{prop:smash-product-connectivity}. Homotopy orbits and small coproducts
of connective spectra are connective. Apply
\Ref{prop:underlying-free-spectral-algebra}.
\end{proof}

\subsection{Connective commutative ring spectra}

\begin{definition}[Connective $\EE_\infty$-ring]
\label{def:connective-e-infinity-ring}
A commutative ring spectrum $A$ is \textbf{connective} if its underlying
spectrum is connective. The full subcategory of connective commutative
ring spectra is denoted
\[
    \CAlg(\Sp)^{\mathrm{cn}}
    \subseteq
    \CAlg(\Sp).
\]
\end{definition}

\begin{proposition}[Presentability, colimits, and limits]
\label{prop:connective-calg-presentable}
The $\infty$-category
\[
    \CAlg(\Sp)^{\mathrm{cn}}
\]
is presentable. The inclusion
\[
    \iota
    \colon
    \CAlg(\Sp)^{\mathrm{cn}}
    \hookrightarrow
    \CAlg(\Sp)
\]
preserves small colimits and admits a right adjoint
\[
    \tau_{\geq0}^{\CAlg}
    \colon
    \CAlg(\Sp)
    \longrightarrow
    \CAlg(\Sp)^{\mathrm{cn}}.
\]
The underlying spectrum of
\[
    \tau_{\geq0}^{\CAlg}(A)
\]
is canonically equivalent to $\tau_{\geq0}U(A)$.

Consequently, for every small diagram
\[
    D
    \colon
    K
    \longrightarrow
    \CAlg(\Sp)^{\mathrm{cn}},
\]
its limit in connective commutative ring spectra is given by
\[
    \varprojlim_K^{\CAlg(\Sp)^{\mathrm{cn}}}D
    \simeq
    \tau_{\geq0}^{\CAlg}
    \left(
        \varprojlim_K^{\CAlg(\Sp)}
        \iota D
    \right).
\]
In general, the inclusion $\iota$ does not preserve finite limits.
\end{proposition}

\begin{proof}
The connective-cover theorem for structured ring spectra identifies the
connective commutative ring spectra as a presentable full subcategory of
$\CAlg(\Sp)$ which is closed under small colimits. Consequently the
inclusion $\iota$ preserves small colimits. Since both categories are
presentable, the adjoint functor theorem supplies a right adjoint
\[
    \tau_{\geq0}^{\CAlg}
    \dashv
    \iota.
\]
The same connective-cover theorem identifies its underlying spectrum:
\[
    U\bigl(\tau_{\geq0}^{\CAlg}(A)\bigr)
    \simeq
    \tau_{\geq0}U(A).
\]
No $t$-structure on the nonstable category $\CAlg(\Sp)$ is being
asserted.

Let $D:K\to\CAlg(\Sp)^{\mathrm{cn}}$ be a small diagram and put
\[
    L
    :=
    \varprojlim_K^{\CAlg(\Sp)}
    \iota D.
\]
For every connective commutative ring spectrum $C$, there are natural
equivalences
\[
\begin{aligned}
    \Map_{\CAlg(\Sp)^{\mathrm{cn}}}
    \bigl(
        C,
        \tau_{\geq0}^{\CAlg}L
    \bigr)
    &\simeq
    \Map_{\CAlg(\Sp)}
    \bigl(
        \iota C,
        L
    \bigr)
    \\
    &\simeq
    \varprojlim_{k\in K}
    \Map_{\CAlg(\Sp)}
    \bigl(
        \iota C,
        \iota D(k)
    \bigr)
    \\
    &\simeq
    \varprojlim_{k\in K}
    \Map_{\CAlg(\Sp)^{\mathrm{cn}}}
    \bigl(
        C,
        D(k)
    \bigr).
\end{aligned}
\]
Yoneda gives the asserted limit formula.

The inclusion does not preserve finite limits in general. Consider the
pullback
\[
    P
    :=
    H\mathbb Z
    \times_{H(\mathbb Z[x])}
    H\mathbb Z,
\]
where both maps are induced by the inclusion of constant polynomials.
The fibre sequence of underlying spectra
\[
    P
    \longrightarrow
    H\mathbb Z\oplus H\mathbb Z
    \longrightarrow
    H(\mathbb Z[x])
\]
gives
\[
    \pi_{-1}(P)
    \simeq
    \operatorname{coker}
    \left(
        \mathbb Z\oplus\mathbb Z
        \longrightarrow
        \mathbb Z[x]
    \right)
    \simeq
    \mathbb Z[x]/\mathbb Z,
\]
which is nonzero. Thus the pullback computed in $\CAlg(\Sp)$ is not
connective; its pullback in $\CAlg(\Sp)^{\mathrm{cn}}$ is obtained by
applying $\tau_{\geq0}^{\CAlg}$.
See \cite{LurieHigherAlgebra,LurieSAG}.
\end{proof}

\begin{corollary}[Connective pushouts]
\label{cor:connective-calg-pushouts}
If
\[
    B
    \longleftarrow
    A
    \longrightarrow
    C
\]
is a span in $\CAlg(\Sp)^{\mathrm{cn}}$, then the pushout
\[
    B\otimes_A C
\]
computed in $\CAlg(\Sp)$ is connective.
\end{corollary}

\begin{proof}
It is a colimit of connective commutative ring spectra. Apply
\Ref{prop:connective-calg-presentable}.
\end{proof}

\subsection{Discrete rings and classical truncation}

\begin{proposition}[Discrete commutative ring spectra]
\label{prop:discrete-commutative-ring-spectra}
The Eilenberg--Mac Lane construction refines to a fully faithful functor
\[
    H
    \colon
    \operatorname{CRing}
    \hookrightarrow
    \CAlg(\Sp)^{\mathrm{cn}}.
\]
Its essential image consists of the commutative ring spectra whose
underlying spectra are discrete.
\end{proposition}

\begin{proof}
The equivalence between discrete spectra and abelian groups of
\Ref{prop:eilenberg-mac-lane-heart} is compatible with finite products.
Commutative algebra objects in the heart are therefore ordinary
commutative rings. Full faithfulness and the description of the essential
image follow.
\end{proof}

\begin{proposition}[Classical truncation adjunction]
\label{prop:pi-zero-h-adjunction}
The functor
\[
    \pi_0
    \colon
    \CAlg(\Sp)^{\mathrm{cn}}
    \longrightarrow
    \operatorname{CRing}
\]
is left adjoint to the Eilenberg--Mac Lane functor. Thus there are natural
equivalences
\[
    \Map_{\operatorname{CRing}}
    \bigl(
        \pi_0(A),R
    \bigr)
    \simeq
    \Map_{\CAlg(\Sp)}
    \bigl(
        A,H R
    \bigr).
\]
The unit of the adjunction is a natural morphism
\[
    A
    \longrightarrow
    H\pi_0(A).
\]
\end{proposition}

\begin{proof}
A morphism from a connective commutative ring spectrum to a discrete
commutative ring spectrum is determined by its effect on degree-zero
homotopy. The multiplicative coherences are exactly those of the induced
ordinary ring homomorphism. See \cite{LurieHigherAlgebra,LurieSAG}.
\end{proof}

\subsection{Truncated commutative ring spectra}

\begin{proposition}[Operadic Postnikov truncations]
\label{prop:operadic-postnikov-truncations}
For every integer $n\geq0$, the truncation functor on spectra lifts to a
functor
\[
    \tau_{\leq n}
    \colon
    \CAlg(\Sp)^{\mathrm{cn}}
    \longrightarrow
    \CAlg(\Sp)^{\mathrm{cn}}
\]
whose underlying spectrum is the ordinary Postnikov truncation. For every
connective commutative ring spectrum $A$, there is a natural equivalence
\[
    A
    \simeq
    \varprojlim_n\tau_{\leq n}A
\]
in $\CAlg(\Sp)$.
\end{proposition}

\begin{proof}
The full subcategory of $n$-truncated objects in commutative ring spectra
is reflective, and the forgetful functor identifies the reflector on
connective objects with the spectral Postnikov truncation. Since the
forgetful functor creates limits, Postnikov convergence in spectra from
\Ref{prop:postnikov-completeness-detection} implies convergence in
commutative ring spectra. See \cite{LurieHigherAlgebra,LurieDAGIV}.
\end{proof}

\section{Affine spectral schemes and classical truncation}
\label{sec:affine-spectral-schemes}

\begin{definition}[Affine spectral scheme]
\label{def:affine-spectral-scheme}
The $\infty$-category of \textbf{affine spectral schemes} is
\[
    \Aff_{\Sp}
    :=
    \bigl(\CAlg(\Sp)^{\mathrm{cn}}\bigr)^{\op}.
\]
For
\[
    A\in\CAlg(\Sp)^{\mathrm{cn}},
\]
the corresponding affine spectral scheme is denoted
\[
    \Spec(A).
\]
\end{definition}

\begin{proposition}[Mapping spaces of affine spectral schemes]
\label{prop:affine-spectral-mapping-spaces}
For connective commutative ring spectra $A$ and $B$, there is a natural
equivalence
\[
    \Map_{\Aff_{\Sp}}
    \bigl(
        \Spec(B),
        \Spec(A)
    \bigr)
    \simeq
    \Map_{\CAlg(\Sp)}(A,B).
\]
\end{proposition}

\begin{proof}
This is the defining property of the opposite $\infty$-category.
\end{proof}

Thus a morphism of affine spectral schemes is represented
contravariantly:
\[
\begin{tikzcd}[column sep=huge]
    A
        \arrow[r, "f"]
    &
    B
    \\
    \Spec(A)
    &
    \Spec(B)
        \arrow[l, "\Spec(f)"'].
\end{tikzcd}
\]

\begin{definition}[Spectral prestacks]
\label{def:spectral-prestacks-affine-stage}
The $\infty$-category of space-valued presheaves on affine spectral
schemes is
\[
    \PSh(\Aff_{\Sp})
    :=
    \Fun
    \bigl(
        (\Aff_{\Sp})^{\op},
        \mathcal S
    \bigr).
\]
Its objects are called \textbf{spectral prestacks}.
\end{definition}

\begin{definition}[Functor of points]
\label{def:spectral-functor-of-points}
For $A\in\CAlg(\Sp)^{\mathrm{cn}}$, the functor of points of
$\Spec(A)$ is
\[
    h_A
    \colon
    (\Aff_{\Sp})^{\op}
    \longrightarrow
    \mathcal S,
\]
defined by
\[
    h_A(\Spec(B))
    :=
    \Map_{\CAlg(\Sp)}(A,B).
\]
\end{definition}

\begin{proposition}[Spectral Yoneda embedding]
\label{prop:spectral-yoneda-embedding}
The assignment
\[
    \Spec(A)
    \longmapsto
    h_A
\]
defines a fully faithful functor
\[
    h
    \colon
    \Aff_{\Sp}
    \hookrightarrow
    \PSh(\Aff_{\Sp}).
\]
\end{proposition}

\begin{proof}
This is the $\infty$-categorical Yoneda lemma.
\end{proof}

\subsection{Finite limits of affine spectral schemes}

\begin{proposition}[Affine spectral fibre products]
\label{prop:affine-spectral-fibre-products}
For a diagram of affine spectral schemes
\[
    \Spec(B)
    \longrightarrow
    \Spec(A)
    \longleftarrow
    \Spec(C),
\]
there is a canonical equivalence
\[
    \Spec(B)
    \times_{\Spec(A)}
    \Spec(C)
    \simeq
    \Spec(B\otimes_A C).
\]
The commutative ring spectrum $B\otimes_A C$ is connective.
\end{proposition}

\begin{proof}
The fibre product in the opposite category is represented by the pushout
of commutative ring spectra. Connectivity follows from
\Ref{cor:connective-calg-pushouts}.
\end{proof}

\begin{corollary}[Finite-limit closure]
\label{cor:affine-spectral-finite-limits}
The $\infty$-category $\Aff_{\Sp}$ admits finite limits.
\end{corollary}

\begin{proof}
The terminal affine spectral scheme is $\Spec(\mathbb S)$, and binary
fibre products exist by \Ref{prop:affine-spectral-fibre-products}.
\end{proof}

\subsection{Classical truncation of affines}

\begin{definition}[Classical truncation of an affine]
\label{def:classical-truncation-affine}
For an affine spectral scheme
\[
    X=\Spec(A),
\]
its \textbf{classical truncation} is the ordinary affine scheme
\[
    X_{\mathrm{cl}}
    :=
    \Spec\bigl(\pi_0(A)\bigr).
\]
The canonical algebra morphism
\[
    A
    \longrightarrow
    H\pi_0(A)
\]
induces a morphism of spectral prestacks
\[
    X_{\mathrm{cl}}
    \longrightarrow
    X.
\]
\end{definition}

\begin{definition}[Underlying topological space]
\label{def:underlying-topological-space-affine-spectral}
The underlying topological space of $X=\Spec(A)$ is
\[
    |X|
    :=
    \left|\Spec\pi_0(A)\right|.
\]
\end{definition}

\begin{proposition}[Functoriality of affine truncation]
\label{prop:affine-classical-truncation-functorial}
The assignment
\[
    \Spec(A)
    \longmapsto
    \Spec\pi_0(A)
\]
defines a functor
\[
    (-)_{\mathrm{cl}}
    \colon
    \Aff_{\Sp}
    \longrightarrow
    \operatorname{AffSch}.
\]
The morphisms $X_{\mathrm{cl}}\to X$ are natural in $X$.
\end{proposition}

\begin{proof}
The functor $\pi_0$ is functorial on connective commutative ring spectra.
Passing to opposite categories gives the asserted functor. Naturality is
the naturality of the adjunction unit from
\Ref{prop:pi-zero-h-adjunction}.
\end{proof}

\begin{remark}[Truncation and fibre products]
\label{rem:classical-truncation-not-left-exact}
Although the inclusion
\[
    \CAlg(\Sp)^{\mathrm{cn}}
    \hookrightarrow
    \CAlg(\Sp)
\]
does not preserve finite limits, classical truncation does preserve
affine spectral fibre products. For every span of connective commutative
ring spectra, the canonical comparison morphism
\[
    \pi_0(B)
    \otimes_{\pi_0(A)}
    \pi_0(C)
    \longrightarrow
    \pi_0(B\otimes_A C)
\]
is an isomorphism. This follows because
\[
    \pi_0
    \colon
    \CAlg(\Sp)^{\mathrm{cn}}
    \longrightarrow
    \operatorname{CRing}
\]
is a left adjoint and therefore preserves pushouts. The corresponding
Cartesian statement for affine schemes is recorded in
\Ref{prop:flat-base-change-classical-truncation}.
\end{remark}

\section{Module spectra and relative commutative algebras}
\label{sec:module-spectra-relative-algebras}

Fix a connective commutative ring spectrum
\[
    A\in\CAlg(\Sp)^{\mathrm{cn}}.
\]

\begin{definition}[Module spectrum]
\label{def:module-spectrum-over-a}
The $\infty$-category of $A$-module spectra is denoted
\[
    \Mod_A(\Sp)
\]
or, when no ambiguity is possible, simply
\[
    \Mod_A.
\]
\end{definition}

\begin{theorem}[Spectral module category]
\label{thm:spectral-module-category}
The category
\[
    \Mod_A
\]
is a presentably stable symmetric monoidal $\infty$-category. Its tensor
product is the relative tensor product
\[
    -\otimes_A-,
\]
its tensor unit is the regular module $A$, and its tensor product preserves
small colimits separately in each variable.
\end{theorem}

\begin{proof}
This is the operadic module construction for a commutative algebra object
in a presentably symmetric monoidal stable $\infty$-category; see
\cite[Theorem~4.5.2.1]{LurieHigherAlgebra}. Limits and colimits are created
in spectra, so stability follows from stability of $\Sp$.
\end{proof}

\begin{construction}[Relative tensor product]
\label{con:spectral-relative-tensor-product}
For a right $A$-module $M$ and a left $A$-module $N$, the relative tensor
product is represented by the geometric realization of the two-sided bar
construction
\[
    M\otimes_A N
    \simeq
    \left|
        [n]
        \longmapsto
        M\otimes A^{\otimes n}\otimes N
    \right|.
\]
Since $A$ is commutative, its left and right module theories are
canonically identified.
\end{construction}

\begin{proposition}[Closedness of module spectra]
\label{prop:module-spectra-closedness}
For every $A$-module $M$, the functor
\[
    -\otimes_A M
    \colon
    \Mod_A
    \longrightarrow
    \Mod_A
\]
admits a right adjoint
\[
    \HHom_A(M,-).
\]
\end{proposition}

\begin{proof}
The functor preserves small colimits by
\Ref{thm:spectral-module-category}. Presentability of $\Mod_A$ and the
adjoint functor theorem supply the right adjoint.
\end{proof}

\subsection{Extension and restriction of scalars}

\begin{definition}[Extension and restriction of scalars]
\label{def:spectral-extension-restriction-scalars}
Let
\[
    f
    \colon
    A
    \longrightarrow
    B
\]
be a morphism of commutative ring spectra. The extension-of-scalars
functor is
\[
    f^*
    :=
    -\otimes_A B
    \colon
    \Mod_A
    \longrightarrow
    \Mod_B.
\]
Its right adjoint is the restriction-of-scalars functor
\[
    f_*
    \colon
    \Mod_B
    \longrightarrow
    \Mod_A.
\]
\end{definition}

The adjunction is displayed by
\[
\begin{tikzcd}[column sep=large]
    \Mod_A
        \arrow[r, shift left=1.1ex, "-\otimes_A B"]
    &
    \Mod_B
        \arrow[l, shift left=1.1ex, "f_*"'].
\end{tikzcd}
\]

\begin{proposition}[Monoidality of spectral base change]
\label{prop:spectral-base-change-monoidal}
The functor
\[
    f^*
    =
    -\otimes_A B
\]
is colimit-preserving and strong symmetric monoidal. In particular, there
is a natural equivalence
\[
    (M\otimes_A N)\otimes_A B
    \simeq
    (M\otimes_A B)
    \otimes_B
    (N\otimes_A B).
\]
\end{proposition}

\begin{proof}
The coCartesian operadic module construction identifies the transition
functor associated to $f$ with extension of scalars and equips it with its
strong symmetric monoidal structure; see
\cite[Theorem~4.5.3.1 and Remark~4.5.3.2]{LurieHigherAlgebra}.
\end{proof}

The coherent monoidality is represented by
\[
\begin{tikzcd}[column sep=large, row sep=large]
    \Mod_A\times\Mod_A
        \arrow[r, "-\otimes_A-"]
        \arrow[d, "f^*\times f^*"']
    &
    \Mod_A
        \arrow[d, "f^*"]
    \\
    \Mod_B\times\Mod_B
        \arrow[r, "-\otimes_B-"']
    &
    \Mod_B.
\end{tikzcd}
\]

\begin{proposition}[Functoriality of module categories]
\label{prop:spectral-module-functoriality}
The assignments
\[
    A
    \longmapsto
    \Mod_A
\]
and
\[
    (A\to B)
    \longmapsto
    (-\otimes_A B)
\]
define a functor
\[
    \mathbf{Mod}_{\Sp}
    \colon
    \CAlg(\Sp)
    \longrightarrow
    \CAlg(\Pr^L_{\mathrm{st}}).
\]
For composable maps $A\to B\to C$, there is a canonical coherent
equivalence
\[
    (-\otimes_A B)\otimes_B C
    \simeq
    -\otimes_A C.
\]
\end{proposition}

\begin{proof}
Straightening the coCartesian family of module categories over
$\CAlg(\Sp)$ gives the asserted functor. The coCartesian composition law
supplies the coherent comparison for composable algebra maps; see
\cite{LurieHigherAlgebra}.
\end{proof}

\subsection{Relative commutative algebra objects}

\begin{proposition}[Relative commutative spectral algebras]
\label{prop:relative-commutative-spectral-algebras}
There is a natural equivalence
\[
    \CAlg(\Mod_A)
    \simeq
    \CAlg(\Sp)_{A/}.
\]
Under this equivalence, a commutative algebra object $B$ in $A$-modules is
identified with its structural morphism
\[
    A
    \longrightarrow
    B.
\]
\end{proposition}

\begin{proof}
This is the commutative specialization of the operadic equivalence between
algebras in a module category and algebras under the base algebra; see
\cite[Corollary~3.4.1.7]{LurieHigherAlgebra}.
\end{proof}

\begin{proposition}[Connective relative algebras]
\label{prop:connective-relative-algebras}
The equivalence of \Ref{prop:relative-commutative-spectral-algebras}
restricts to an equivalence
\[
    \CAlg(\Mod_A)^{\mathrm{cn}}
    \simeq
    \CAlg(\Sp)^{\mathrm{cn}}_{A/},
\]
where connectivity on the left is detected on the underlying $A$-module.
Consequently,
\[
    \Aff_{\Sp}/\Spec(A)
    \simeq
    \bigl(
        \CAlg(\Mod_A)^{\mathrm{cn}}
    \bigr)^{\op}.
\]
\end{proposition}

\begin{proof}
The equivalence preserves underlying spectra and therefore preserves
connective objects. Passing to opposite categories gives the affine
statement.
\end{proof}

\subsection{The induced module \texorpdfstring{$t$}{t}-structure}

\begin{definition}[Module Postnikov $t$-structure]
\label{def:module-postnikov-t-structure}
Let
\[
    \Mod_{A,\geq0}
\]
be the full subcategory of $A$-modules whose underlying spectra are
connective, and let
\[
    \Mod_{A,\leq0}
\]
be the full subcategory whose underlying spectra are coconnective.
\end{definition}

\begin{proposition}[Induced module $t$-structure]
\label{prop:induced-module-t-structure}
The pair
\[
    (\Mod_{A,\geq0},\Mod_{A,\leq0})
\]
is an accessible, left complete, and right complete $t$-structure on
$\Mod_A$. The forgetful functor
\[
    U_A
    \colon
    \Mod_A
    \longrightarrow
    \Sp
\]
is $t$-exact and creates Postnikov truncations.
\end{proposition}

\begin{proof}
Limits and colimits of $A$-modules are created in spectra. The module
structure on a truncation is induced by the connective multiplication of
$A$, and the orthogonality and truncation axioms are detected after
forgetting to spectra. Completeness follows from completeness of the
Postnikov $t$-structure and creation of limits and colimits. See
\cite{LurieHigherAlgebra,LurieSAG}.
\end{proof}

\begin{proposition}[Heart of the module category]
\label{prop:heart-module-category}
There is a natural equivalence of abelian categories
\[
    \Mod_A^\heartsuit
    \simeq
    \operatorname{Mod}_{\pi_0(A)}.
\]
\end{proposition}

\begin{proof}
A discrete $A$-module is determined by its degree-zero homotopy group. The
$A$-action factors uniquely through the discrete commutative ring spectrum
$H\pi_0(A)$, and hence is equivalent to an ordinary $\pi_0(A)$-module
structure.
\end{proof}

\begin{proposition}[Exactness and right $t$-exactness of base change]
\label{prop:base-change-exact-right-t-exact}
Let $f:A\to B$ be a morphism of connective commutative ring spectra.
Then:
\begin{enumerate}[label=(\roman*)]
    \item the extension-of-scalars functor
    \[
        f^*
        \colon
        \Mod_A
        \longrightarrow
        \Mod_B
    \]
    is exact as a functor of stable $\infty$-categories;

    \item the functor $f^*$ is right $t$-exact;

    \item the restriction-of-scalars functor $f_*$ is $t$-exact.
\end{enumerate}
\end{proposition}

\begin{proof}
The functor $f^*$ preserves all colimits. A colimit-preserving functor
between stable $\infty$-categories preserves finite colimits and therefore
finite limits, proving exactness. Since $B$ is connective as an
$A$-module, the relative tensor product of a connective $A$-module with
$B$ is connective, proving right $t$-exactness. Restriction of scalars does
not change the underlying spectrum and is therefore $t$-exact.
\end{proof}

\begin{remark}[Stable exactness is not flatness]
\label{rem:stable-exactness-not-flatness}
Assertion (i) of \Ref{prop:base-change-exact-right-t-exact} holds for every
morphism of commutative ring spectra and carries no geometric flatness
content. In this chapter the word \emph{flat} is reserved for the
$t$-structural condition introduced in
\Ref{def:flat-module-spectrum}.
\end{remark}

\section{Spectral flatness and finiteness conditions}
\label{sec:spectral-flatness-finiteness}

The canonical $t$-structures on module categories distinguish geometric
flatness from the automatic exactness of stable base change.

\subsection{Flat module spectra}

\begin{definition}[Flat module spectrum]
\label{def:flat-module-spectrum}
Let $A$ be a connective commutative ring spectrum. A connective
$A$-module $M$ is \textbf{flat} if the functor
\[
    -\otimes_A M
    \colon
    \Mod_A
    \longrightarrow
    \Mod_A
\]
is $t$-exact.
\end{definition}

The right $t$-exactness is automatic for a connective module. Flatness is
therefore the additional assertion that tensoring with $M$ preserves
coconnective modules.

\begin{theorem}[Spectral Lazard theorem]
\label{thm:spectral-lazard}
Let $A$ be a connective commutative ring spectrum and let
$M\in\Mod_{A,\geq0}$. The following conditions are equivalent.
\begin{enumerate}[label=(\roman*)]
    \item The module $M$ is flat.

    \item The ordinary $\pi_0(A)$-module $\pi_0(M)$ is flat and, for every
    integer $n$, the canonical morphism
    \[
        \pi_n(A)
        \otimes_{\pi_0(A)}
        \pi_0(M)
        \longrightarrow
        \pi_n(M)
    \]
    is an isomorphism.

    \item For every $A$-module $N$ and every integer $n$, the canonical
    morphism
    \[
        \pi_n(N)
        \otimes_{\pi_0(A)}
        \pi_0(M)
        \longrightarrow
        \pi_n(N\otimes_A M)
    \]
    is an isomorphism.

    \item The module $M$ is a filtered colimit of finite free
    $A$-modules.
\end{enumerate}
\end{theorem}

\begin{proof}
The equivalence of the $t$-exactness, homotopy-group, and filtered-colimit
conditions is the spectral form of Lazard's theorem; see
\cite{LurieHigherAlgebra,LurieSAG}. The implication from (iii) to (ii) is
obtained by taking $N=A$. Condition (ii), together with the Tor spectral
sequence for $N\otimes_A M$, forces all higher Tor contributions to
vanish and yields (iii). Condition (iv) implies (iii) because tensor
products and homotopy groups commute with filtered colimits and the formula
is immediate for finite free modules. The converse is the spectral
flat-module classification.
\end{proof}

\begin{corollary}[Homotopy groups of a flat module]
\label{cor:flat-modules-detected-homotopy}
If $M$ is a flat $A$-module, then its graded homotopy groups are obtained
from $\pi_0(M)$ by base change:
\[
    \pi_n(M)
    \simeq
    \pi_n(A)
    \otimes_{\pi_0(A)}
    \pi_0(M)
\]
for every integer $n$.
\end{corollary}

\begin{proof}
This is condition (ii) of \Ref{thm:spectral-lazard}.
\end{proof}

\begin{proposition}[Permanence of flat modules]
\label{prop:flat-module-permanence}
Let $A$ be a connective commutative ring spectrum.
\begin{enumerate}[label=(\roman*)]
    \item Finite free and finite projective $A$-modules are flat.

    \item Filtered colimits of flat $A$-modules are flat.

    \item If $M$ and $N$ are flat $A$-modules, then
    \[
        M\otimes_A N
    \]
    is flat over $A$.

    \item Let $A\to B$ be a flat morphism of connective commutative ring
    spectra. If $N$ is flat as a $B$-module, then $N$, regarded by
    restriction of scalars as an $A$-module, is flat over $A$.

    \item Let $A\to B$ be any morphism of connective commutative ring
    spectra. If $M$ is flat over $A$, then
    \[
        M\otimes_A B
    \]
    is flat over $B$.
\end{enumerate}
\end{proposition}

\begin{proof}
A finite free module is a finite direct sum of copies of $A$, so tensoring
with it is a finite direct sum of the identity functor and is therefore
$t$-exact. If $M$ is a retract of a flat module $F$, then
\[
    -\otimes_A M
\]
is a retract of the $t$-exact functor $-\otimes_A F$. Both the connective
and coconnective subcategories are closed under retracts, so $M$ is flat.
This proves (i).

For a filtered diagram of flat modules $\{M_\alpha\}$, there is a natural
equivalence
\[
    K\otimes_A
    \left(
        \varinjlim_\alpha M_\alpha
    \right)
    \simeq
    \varinjlim_\alpha
    (K\otimes_A M_\alpha).
\]
Filtered colimits preserve both connective and coconnective objects in
the Postnikov $t$-structure, because homotopy groups commute with filtered
colimits. Thus the displayed tensor functor is $t$-exact, proving (ii).

For (iii), associativity gives
\[
    -\otimes_A(M\otimes_A N)
    \simeq
    \bigl(
        -\otimes_A M
    \bigr)
    \otimes_A N.
\]
The right side is a composite of $t$-exact functors.

For (iv), let $K$ be an $A$-module. There is a natural equivalence of
underlying $A$-modules
\[
    K\otimes_A N
    \simeq
    (K\otimes_A B)\otimes_B N.
\]
Extension of scalars from $A$ to $B$ is $t$-exact by flatness of
$A\to B$, tensoring with $N$ over $B$ is $t$-exact by flatness of $N$,
and restriction of scalars from $B$ to $A$ is $t$-exact. Their composite
is therefore $t$-exact.

For (v), let $K$ be a $B$-module. Then
\[
\begin{aligned}
    K\otimes_B(M\otimes_A B)
    &\simeq
    K\otimes_A M,
\end{aligned}
\]
where $K$ on the right is regarded as an $A$-module. Restriction from
$B$ to $A$ does not change the underlying spectrum and is $t$-exact.
Since $M$ is flat over $A$, the right side depends $t$-exactly on $K$.
Thus $M\otimes_A B$ is flat over $B$.
\end{proof}

\subsection{Flat morphisms of connective commutative ring spectra}

\begin{definition}[Flat spectral algebra morphism]
\label{def:flat-spectral-algebra-morphism}
A morphism
\[
    f
    \colon
    A
    \longrightarrow
    B
\]
in $\CAlg(\Sp)^{\mathrm{cn}}$ is \textbf{flat} if $B$, regarded as an
$A$-module, is flat.
\end{definition}

\begin{proposition}[Homotopy groups of a flat algebra]
\label{prop:homotopy-groups-flat-algebra}
If $A\to B$ is flat, then $\pi_0(B)$ is flat over $\pi_0(A)$ and there are
natural isomorphisms
\[
    \pi_n(B)
    \simeq
    \pi_n(A)
    \otimes_{\pi_0(A)}
    \pi_0(B)
\]
for every integer $n$.
\end{proposition}

\begin{proof}
Apply condition (ii) of \Ref{thm:spectral-lazard} to the $A$-module $B$.
\end{proof}

\begin{proposition}[Permanence of flat algebra morphisms]
\label{prop:flat-algebra-permanence}
Flat morphisms in $\CAlg(\Sp)^{\mathrm{cn}}$ are stable under equivalence,
composition, base change, and filtered colimits in the arrow category.
\end{proposition}\begin{proof}
Stability under equivalence is immediate.

Let
\[
    A
    \longrightarrow
    B
    \longrightarrow
    C
\]
be flat morphisms. Extension of scalars from $A$ to $C$ is canonically
equivalent to the composite
\[
    \Mod_A
    \xrightarrow{-\otimes_AB}
    \Mod_B
    \xrightarrow{-\otimes_BC}
    \Mod_C.
\]
Both functors are $t$-exact, so their composite is $t$-exact. Hence
$A\to C$ is flat.

For base change, let $A\to A'$ be any morphism and put
\[
    B'
    :=
    B\otimes_AA'.
\]
For every $A'$-module $M$, there is a canonical equivalence of underlying
spectra
\[
    M\otimes_{A'}B'
    \simeq
    M\otimes_AB,
\]
where $M$ on the right is restricted to an $A$-module. Restriction of
scalars from $A'$ to $A$ is $t$-exact, and tensoring with $B$ over $A$ is
$t$-exact by flatness. Thus
\[
    -\otimes_{A'}B'
\]
is $t$-exact, and $A'\to B'$ is flat.

Finally, let
\[
    \{A_\alpha\to B_\alpha\}_{\alpha}
\]
be a filtered diagram of flat morphisms and put
\[
    A:=\varinjlim_\alpha A_\alpha,
    \qquad
    B:=\varinjlim_\alpha B_\alpha.
\]
Filtered colimits commute with homotopy groups. The flat homotopy-group
formula at each stage gives
\[
    \pi_n(B_\alpha)
    \simeq
    \pi_n(A_\alpha)
    \otimes_{\pi_0(A_\alpha)}
    \pi_0(B_\alpha).
\]
Passing to filtered colimits yields
\[
    \pi_n(B)
    \simeq
    \pi_n(A)
    \otimes_{\pi_0(A)}
    \pi_0(B).
\]
Moreover, $\pi_0(B)$ is flat over $\pi_0(A)$, since flatness of ordinary
modules is preserved under filtered colimits in the arrow category.
The flatness criterion of
\Ref{thm:spectral-lazard} therefore implies that $A\to B$ is flat.
\end{proof}

\begin{proposition}[Classical truncation and affine fibre products]
\label{prop:flat-base-change-classical-truncation}
Let
\[
\begin{tikzcd}[column sep=large, row sep=large]
    A
        \arrow[r]
        \arrow[d]
    &
    B
        \arrow[d]
    \\
    A'
        \arrow[r]
    &
    B':=B\otimes_AA'
\end{tikzcd}
\]
be a pushout square in connective commutative ring spectra. Then the
canonical morphism
\[
    \pi_0(B)
    \otimes_{\pi_0(A)}
    \pi_0(A')
    \longrightarrow
    \pi_0(B')
\]
is an isomorphism. Equivalently, the square of affine schemes
\[
\begin{tikzcd}[column sep=large, row sep=large]
    \Spec\pi_0(B')
        \arrow[r]
        \arrow[d]
    &
    \Spec\pi_0(B)
        \arrow[d]
    \\
    \Spec\pi_0(A')
        \arrow[r]
    &
    \Spec\pi_0(A)
\end{tikzcd}
\]
is Cartesian.
\end{proposition}

\begin{proof}
The adjunction of \Ref{prop:pi-zero-h-adjunction} shows that
\[
    \pi_0
    \colon
    \CAlg(\Sp)^{\mathrm{cn}}
    \longrightarrow
    \operatorname{CRing}
\]
is a left adjoint. It therefore preserves all small colimits, in
particular pushouts. Hence
\[
    \pi_0(B')
    \simeq
    \pi_0(B)
    \otimes_{\pi_0(A)}
    \pi_0(A').
\]
Passing to opposite categories gives the affine statement.
\end{proof}

\subsection{Finite presentation and finite cell algebras}

\begin{definition}[Finite presentation]
\label{def:finitely-presented-spectral-algebra}
A morphism
\[
    A
    \longrightarrow
    B
\]
of connective commutative ring spectra is \textbf{of finite presentation}
if $B$ is compact as an object of the $\infty$-category
\[
    \CAlg(\Sp)_{A/}.
\]
A morphism of affine spectral schemes is of finite presentation if its
contravariant algebra morphism is of finite presentation. At the
algebraic level this compactness condition is also commonly called
``locally of finite presentation''; for an affine geometric morphism the
global quasi-compactness and separatedness conditions are automatic.
\end{definition}

\begin{definition}[Finite cell algebra]
\label{def:finite-cell-spectral-algebra}
A commutative $A$-algebra is a \textbf{finite cell algebra} if it is
obtained from $A$ by finitely many pushouts of free algebra maps induced by
maps between finite spectra.
\end{definition}

\begin{proposition}[Compactness criterion]
\label{prop:finite-presentation-compactness-criterion}
A connective commutative $A$-algebra is of finite presentation if and only
if it is a retract of a finite cell commutative $A$-algebra.
\end{proposition}

\begin{proof}
The $\infty$-category of commutative $A$-algebras is compactly generated by
free algebras on finite spectra. Compact objects are therefore retracts of
finite cell objects; see \cite{LurieHigherAlgebra,LurieSAG}.
\end{proof}

\begin{proposition}[Permanence of finite presentation]
\label{prop:finite-presentation-permanence}
Morphisms of finite presentation are stable under equivalence,
composition, and base change.
\end{proposition}

\begin{proof}
Equivalence stability is immediate. Compactness is preserved under the
base-change left adjoint because its right adjoint preserves filtered
colimits. Composition follows by expressing mapping spaces out of the
composite algebra through the two compact stages. Equivalently, one may
use finite cell presentations and concatenate or base-change the finite
cell attachments.
\end{proof}

\begin{lemma}[Compact arrows in an undercategory]
\label{lem:compact-arrow-undercategory}
Let $\mathcal C$ be a presentable $\infty$-category and let
\[
    P
    \longrightarrow
    B
\]
be a morphism between compact objects of $\mathcal C$. Then $B$, regarded
as an object of the undercategory $\mathcal C_{P/}$, is compact.
\end{lemma}

\begin{proof}
Let $\{C_\alpha\}_\alpha$ be a filtered diagram in $\mathcal C_{P/}$.
The mapping space
\[
    \Map_{\mathcal C_{P/}}
    \left(
        B,
        \varinjlim_\alpha C_\alpha
    \right)
\]
is the homotopy fibre, over the structural map
$P\to\varinjlim_\alpha C_\alpha$, of
\[
    \Map_{\mathcal C}
    \left(
        B,
        \varinjlim_\alpha C_\alpha
    \right)
    \longrightarrow
    \Map_{\mathcal C}
    \left(
        P,
        \varinjlim_\alpha C_\alpha
    \right).
\]
Compactness of $P$ and $B$ identifies the two mapping spaces with the
corresponding filtered colimits. Filtered colimits of spaces commute with
finite limits, hence with the displayed homotopy fibre. Therefore
\[
    \Map_{\mathcal C_{P/}}
    \left(
        B,
        \varinjlim_\alpha C_\alpha
    \right)
    \simeq
    \varinjlim_\alpha
    \Map_{\mathcal C_{P/}}(B,C_\alpha),
\]
which proves compactness in the undercategory.
\end{proof}

\subsection{Perfect, projective, and locally free modules}

\begin{definition}[Perfect module]
\label{def:perfect-module-spectrum}
An $A$-module $M$ is \textbf{perfect} if it belongs to the thick
subcategory of $\Mod_A$ generated by the tensor unit $A$.
\end{definition}

\begin{proposition}[Perfect equals dualizable]
\label{prop:perfect-equals-dualizable}
An $A$-module is perfect if and only if it is dualizable in the symmetric
monoidal $\infty$-category $\Mod_A$.
\end{proposition}

\begin{proof}
The unit $A$ is compact and dualizable, and the dualizable objects form a
thick subcategory. Conversely, a dualizable $A$-module is compact because
its mapping functor is represented after tensoring with its dual, and the
compact objects of $\Mod_A$ form the thick subcategory generated by $A$.
See \cite{LurieHigherAlgebra}.
\end{proof}

\begin{definition}[Tor-amplitude]
\label{def:tor-amplitude}
An $A$-module $M$ has \textbf{Tor-amplitude in $[a,b]$} if, for every
discrete $A$-module $N$, the module
\[
    M\otimes_A N
\]
has homotopy groups only in degrees $a$ through $b$.
\end{definition}

\begin{definition}[Finite projective module]
\label{def:finite-projective-module-spectrum}
An $A$-module $M$ is \textbf{finite projective} if it is a retract of a
finite free $A$-module.
\end{definition}

\begin{proposition}[Characterization of finite projectives]
\label{prop:finite-projective-characterization}
For an $A$-module $M$, the following are equivalent.
\begin{enumerate}[label=(\roman*)]
    \item The module $M$ is finite projective.

    \item The module $M$ is perfect and has Tor-amplitude contained in
    $[0,0]$.

    \item The module $M$ is perfect, connective, and flat.
\end{enumerate}
\end{proposition}
\begin{proof}
A finite free \(A\)-module is perfect, connective, flat, and has
Tor-amplitude in \([0,0]\). Each of these properties is stable under
retracts. Thus \textup{(i)} implies both \textup{(ii)} and
\textup{(iii)}.

Assume \textup{(iii)}. Since \(M\) is perfect, its base change
\[
    M_0
    :=
    M\otimes_AH\pi_0(A)
\]
is perfect over the discrete ring spectrum \(H\pi_0(A)\). Flatness gives
\[
    \pi_n(M_0)
    \simeq
    \pi_n(H\pi_0(A))
    \otimes_{\pi_0(A)}
    \pi_0(M),
\]
so \(M_0\) is discrete with
\[
    \pi_0(M_0)
    \simeq
    \pi_0(M).
\]

Since \(M_0\) is perfect and discrete over the discrete ring spectrum
\(H\pi_0(A)\), the ordinary \(\pi_0(A)\)-module
\[
    \pi_0(M_0)
    \simeq
    \pi_0(M)
\]
admits a finite resolution by finitely generated projective modules. In
particular, \(\pi_0(M)\) is finitely presented over \(\pi_0(A)\).

Since \(M\) is flat, \(\pi_0(M)\) is flat over \(\pi_0(A)\). A finitely
presented flat module is finitely generated projective. Hence
\(\pi_0(M)\) is finitely generated projective over \(\pi_0(A)\).

Choose an idempotent matrix
\[
    e
    \in
    M_r(\pi_0(A))
\]
whose image is \(\pi_0(M)\). The same idempotent defines an idempotent
endomorphism of the finite free \(A\)-module \(A^{\oplus r}\). Since
idempotents split in \(\Mod_A\), it determines a finite projective
\(A\)-module \(P\) satisfying
\[
    \pi_0(P)
    \simeq
    \pi_0(M).
\]

Because \(P\) is a retract of a finite free module, the canonical map
\[
    \pi_0
    \Map_{\Mod_A}(P,M)
    \longrightarrow
    \Hom_{\pi_0(A)}
    \bigl(
        \pi_0(P),
        \pi_0(M)
    \bigr)
\]
is surjective. Choose a morphism
\[
    u
    \colon
    P
    \longrightarrow
    M
\]
which induces the chosen isomorphism on \(\pi_0\).

Both \(P\) and \(M\) are flat. Therefore, for every integer \(n\),
\[
    \pi_n(P)
    \simeq
    \pi_n(A)
    \otimes_{\pi_0(A)}
    \pi_0(P)
\]
and
\[
    \pi_n(M)
    \simeq
    \pi_n(A)
    \otimes_{\pi_0(A)}
    \pi_0(M).
\]
Under these identifications, \(\pi_n(u)\) is obtained by tensoring
\(\pi_0(u)\) with \(\pi_n(A)\) and is therefore an isomorphism. Hence
\(u\) is an equivalence and \(M\) is finite projective. This proves
\[
    \textup{(iii)}
    \Longrightarrow
    \textup{(i)}.
\]

Now assume \textup{(ii)}. Put
\[
    M_0
    :=
    M\otimes_AH\pi_0(A).
\]
By the Tor-amplitude condition, \(M_0\) has homotopy only in degree zero.
Since \(M\) is perfect over the connective ring spectrum \(A\), it is
bounded below.

Let
\[
    I
    :=
    \operatorname{fib}
    \bigl(
        A\longrightarrow H\pi_0(A)
    \bigr).
\]
The map \(A\to H\pi_0(A)\) induces an isomorphism on degree-zero
homotopy, so
\[
    I\in\Mod_{A,\geq1}.
\]
Suppose \(M\neq0\), and let \(r\) be the least integer for which
\[
    \pi_r(M)\neq0.
\]
Then
\[
    M\in\Mod_{A,\geq r}.
\]
The bar construction for the relative tensor product shows that, for
\(P\in\Mod_{A,\geq a}\) and \(Q\in\Mod_{A,\geq b}\),
\[
    P\otimes_AQ
    \in
    \Mod_{A,\geq a+b}.
\]
Consequently,
\[
    I\otimes_AM
    \in
    \Mod_{A,\geq r+1}.
\]
Tensoring the fibre sequence
\[
    I
    \longrightarrow
    A
    \longrightarrow
    H\pi_0(A)
\]
with \(M\) gives a fibre sequence
\[
    I\otimes_AM
    \longrightarrow
    M
    \longrightarrow
    M_0.
\]
It follows that
\[
    \pi_r(M)
    \xrightarrow{\simeq}
    \pi_r(M_0).
\]
Since \(M_0\) is discrete, one has \(r=0\). Thus \(M\) is connective.

We next prove that \(M\) is flat. Let
\[
    F
    :=
    -\otimes_AM.
\]
The connectivity of \(M\) implies that \(F\) is right \(t\)-exact. By
the Tor-amplitude condition, \(F\) carries every discrete \(A\)-module
to a discrete spectrum.

Let \(N\in\Mod_{A,\leq0}\). Right completeness gives
\[
    N
    \simeq
    \varinjlim_k\tau_{\geq-k}N.
\]
Each bounded object \(\tau_{\geq-k}N\) is obtained by finitely many
extensions from the shifted discrete modules
\[
    H\pi_j(N)[j],
    \qquad
    -k\leq j\leq0.
\]
The functor \(F\) is exact and carries discrete modules to discrete
spectra, so
\[
    F\bigl(\tau_{\geq-k}N\bigr)
    \in
    \Sp_{\leq0}
\]
for every \(k\). Since filtered colimits of coconnective spectra are
coconnective,
\[
    F(N)
    \simeq
    \varinjlim_kF\bigl(\tau_{\geq-k}N\bigr)
    \in
    \Sp_{\leq0}.
\]
Thus \(F\) is also left \(t\)-exact, and \(M\) is flat. Therefore \(M\)
satisfies \textup{(iii)}, and the implication
\[
    \textup{(iii)}
    \Longrightarrow
    \textup{(i)}
\]
already proved above shows that \(M\) is finite projective. This proves
\[
    \textup{(ii)}
    \Longrightarrow
    \textup{(i)}.
\]
\end{proof}

\begin{definition}[Locally free module]
\label{def:locally-free-affine-module}
Let $A$ be connective. An $A$-module $M$ is \textbf{locally free of rank
$n$} if there exists a Zariski covering family
\[
    \{A\to A_i\}_{i\in I}
\]
as defined in \Ref{def:affine-zariski-cover}, such that
\[
    M\otimes_AA_i
    \simeq
    A_i^{\oplus n}
\]
for every $i$.
\end{definition}

\begin{remark}
\label{rem:local-free-definition-forward-reference}
The preceding definition uses the spectral Zariski topology constructed in
\Ref{sec:affine-opens-zariski-descent}. Until that construction is
complete, finite projective modules are the affine algebraic substitute.
\end{remark}

\section{Spectral prestacks and quasi-coherent modules}
\label{sec:spectral-prestacks-qcoh}

\subsection{The affine quasi-coherent functor}

\begin{definition}[Affine quasi-coherent modules]
\label{def:affine-spectral-qcoh}
For an affine spectral scheme $\Spec(A)$, define
\[
    \QCoh_{\mathrm{aff}}(\Spec(A))
    :=
    \Mod_A.
\]
For a morphism
\[
    \Spec(B)
    \longrightarrow
    \Spec(A),
\]
the pullback functor is extension of scalars
\[
    -\otimes_A B
    \colon
    \Mod_A
    \longrightarrow
    \Mod_B.
\]
\end{definition}

\begin{proposition}[Affine quasi-coherent coefficient system]
\label{prop:affine-qcoh-coefficient-system}
The assignments of \Ref{def:affine-spectral-qcoh} define a functor
\[
    \QCoh_{\mathrm{aff}}
    \colon
    (\Aff_{\Sp})^{\op}
    \longrightarrow
    \CAlg(\Pr^L_{\mathrm{st}}).
\]
\end{proposition}

\begin{proof}
This is the restriction of the functor
$\mathbf{Mod}_{\Sp}$ from
\Ref{prop:spectral-module-functoriality} to connective commutative ring
spectra.
\end{proof}

\subsection{Affine points and right Kan extension}

\begin{definition}[Affine points of a prestack]
\label{def:affine-points-prestack}
For a spectral prestack
\[
    X\in\PSh(\Aff_{\Sp}),
\]
let
\[
    \Aff_{\Sp}/X
\]
denote the $\infty$-category of pairs
\[
    (\Spec(A),x),
    \qquad
    x\in X(A),
\]
equivalently the slice of the Yoneda embedding over $X$.
\end{definition}

\begin{definition}[Quasi-coherent modules on a spectral prestack]
\label{def:qcoh-spectral-prestack}
For a spectral prestack $X$, define
\[
    \QCoh(X)
    :=
    \varprojlim_{(\Spec(A)\to X)
    \in(\Aff_{\Sp}/X)^{\op}}
    \Mod_A.
\]
The limit is formed in
\[
    \CAlg(\Pr^L_{\mathrm{st}}).
\]
\end{definition}

The definition is represented schematically by
\[
\begin{tikzcd}[column sep=large]
    (\Aff_{\Sp}/X)^{\op}
        \arrow[r]
    &
    \CAlg(\Pr^L_{\mathrm{st}})
    \\
    (\Spec(A)\to X)
        \arrow[r, mapsto]
    &
    \Mod_A.
\end{tikzcd}
\]

\begin{theorem}[Affine right Kan extension]
\label{thm:qcoh-affine-right-kan-extension}
Let
\[
    h
    \colon
    \Aff_{\Sp}
    \longrightarrow
    \PSh(\Aff_{\Sp})
\]
be the Yoneda embedding. The construction of
\Ref{def:qcoh-spectral-prestack} defines a functor
\[
    \QCoh
    \colon
    \PSh(\Aff_{\Sp})^{\op}
    \longrightarrow
    \CAlg(\Pr^L_{\mathrm{st}})
\]
which is the right Kan extension of
\[
    \QCoh_{\mathrm{aff}}
    \colon
    (\Aff_{\Sp})^{\op}
    \longrightarrow
    \CAlg(\Pr^L_{\mathrm{st}})
\]
along the opposite Yoneda embedding
\[
    h^{\op}
    \colon
    (\Aff_{\Sp})^{\op}
    \longrightarrow
    \PSh(\Aff_{\Sp})^{\op}.
\]

Equivalently, $\QCoh$ is the unique limit-preserving functor on
$\PSh(\Aff_{\Sp})^{\op}$ whose restriction along
$h^{\op}$ is $\QCoh_{\mathrm{aff}}$.
\end{theorem}

\begin{proof}
For a spectral prestack $X$, the pointwise formula for right Kan extension
along $h^{\op}$ is
\[
    \operatorname{Ran}_{h^{\op}}
    (\QCoh_{\mathrm{aff}})(X)
    \simeq
    \varprojlim_
    {(\Spec(A)\to X)
    \in
    (\Aff_{\Sp}/X)^{\op}}
    \Mod_A.
\]
This is exactly the definition of $\QCoh(X)$ in
\Ref{def:qcoh-spectral-prestack}.

Functoriality follows from functoriality of the affine-point categories:
a morphism $X\to Y$ induces a functor
\[
    \Aff_{\Sp}/X
    \longrightarrow
    \Aff_{\Sp}/Y
\]
by composition. Passing to opposites and applying the universal property
of limits produces the pullback functor
\[
    \QCoh(Y)
    \longrightarrow
    \QCoh(X).
\]

The opposite Yoneda embedding generates
\[
    \PSh(\Aff_{\Sp})^{\op}
\]
under limits. Therefore a limit-preserving extension is uniquely
determined by its restriction to representables. The right Kan extension
is limit-preserving, so it has the claimed universal property.
\end{proof}

\begin{corollary}[Recovery on affine spectral schemes]
\label{cor:qcoh-recovery-affines}
For every connective commutative ring spectrum $A$, there is a canonical
symmetric monoidal equivalence
\[
    \QCoh(\Spec(A))
    \simeq
    \Mod_A.
\]
\end{corollary}\begin{proof}
For the representable prestack
\[
    X=\Spec(A),
\]
the identity map
\[
    \id_X
    \colon
    X
    \longrightarrow
    X
\]
is terminal in the affine-point category
\[
    \Aff_{\Sp}/X.
\]
It is therefore initial in
\[
    (\Aff_{\Sp}/X)^{\op}.
\]
A limit of a diagram indexed by an $\infty$-category with an initial
object is canonically its value at that initial object. Hence
\[
    \QCoh(\Spec(A))
    \simeq
    \Mod_A.
\]
The equivalence is symmetric monoidal because the defining limit is
formed in
\[
    \CAlg(\Pr^L_{\mathrm{st}}).
\]
\end{proof}

\begin{corollary}[Colimits of prestacks]
\label{cor:qcoh-colimits-to-limits}
For every small diagram of spectral prestacks $X_\alpha$, there is a
canonical equivalence
\[
    \QCoh
    \left(
        \varinjlim_\alpha X_\alpha
    \right)
    \simeq
    \varprojlim_\alpha\QCoh(X_\alpha).
\]
\end{corollary}

\begin{proof}
Apply the limit-preservation assertion of
\Ref{thm:qcoh-affine-right-kan-extension} to the opposite functor.
\end{proof}

\subsection{Pullback, direct image, and the structure sheaf}

\begin{proposition}[Quasi-coherent pullback]
\label{prop:qcoh-pullback}
Every morphism of spectral prestacks
\[
    f
    \colon
    X
    \longrightarrow
    Y
\]
induces a colimit-preserving strong symmetric monoidal functor
\[
    f^*
    \colon
    \QCoh(Y)
    \longrightarrow
    \QCoh(X).
\]
For affine spectral schemes it agrees with extension of scalars.
\end{proposition}

\begin{proof}
Functoriality of the right Kan extension gives the transition functor.
The target of the right Kan extension is
$\CAlg(\Pr^L_{\mathrm{st}})$, so the transition is colimit-preserving and
strong symmetric monoidal. The affine identification follows from
\Ref{cor:qcoh-recovery-affines} and the defining affine coefficient
system.
\end{proof}

\begin{corollary}[Quasi-coherent direct image]
\label{cor:qcoh-direct-image}
The pullback functor of \Ref{prop:qcoh-pullback} admits a right adjoint
\[
    f_*
    \colon
    \QCoh(X)
    \longrightarrow
    \QCoh(Y).
\]
\end{corollary}

\begin{proof}
The functor $f^*$ preserves small colimits between presentable
$\infty$-categories. The adjoint functor theorem supplies a right adjoint.
\end{proof}

\begin{definition}[Structure sheaf]
\label{def:structure-sheaf-prestack}
The \textbf{structure sheaf} of a spectral prestack $X$ is the tensor unit
\[
    \mathcal O_X
    \in
    \QCoh(X).
\]
\end{definition}

\begin{definition}[Global sections]
\label{def:global-sections-qcoh}
For $\mathcal F\in\QCoh(X)$, define its spectrum of global sections by
\[
    \Gamma(X,\mathcal F)
    :=
    \map_{\QCoh(X)}
    (\mathcal O_X,\mathcal F).
\]
\end{definition}

\begin{proposition}[Pullback of the structure sheaf]
\label{prop:pullback-structure-sheaf}
For every morphism $f:X\to Y$, there is a canonical equivalence
\[
    f^*\mathcal O_Y
    \simeq
    \mathcal O_X.
\]
\end{proposition}

\begin{proof}
The functor $f^*$ is strong symmetric monoidal and therefore preserves the
tensor unit.
\end{proof}

\subsection{Connectivity on prestacks}

\begin{definition}[Objectwise connective quasi-coherent module]
\label{def:objectwise-connective-qcoh}
A quasi-coherent module $\mathcal F\in\QCoh(X)$ is
\textbf{objectwise connective} if, for every affine point
\[
    x
    \colon
    \Spec(A)
    \longrightarrow
    X,
\]
its pullback
\[
    x^*\mathcal F
\]
belongs to $\Mod_{A,\geq0}$. Objectwise coconnective modules are defined
similarly.
\end{definition}

\begin{remark}
\label{rem:prestack-t-structure}
For arbitrary prestacks, the objectwise connective and coconnective
subcategories need not have all the locality properties required in
geometry. After Zariski descent is established, they define the canonical
$t$-structure on quasi-coherent modules over spectral schemes in
\Ref{thm:qcoh-spectral-scheme-t-structure}.
\end{remark}

\section{Affine open immersions and spectral Zariski descent}
\label{sec:affine-opens-zariski-descent}

\subsection{Affine monomorphisms and open immersions}

\begin{proposition}[Affine monomorphism criterion]
\label{prop:affine-monomorphism-criterion}
Let $A\to B$ be a morphism of connective commutative ring spectra. The
following are equivalent.
\begin{enumerate}[label=(\roman*)]
    \item The morphism
    \[
        \Spec(B)
        \longrightarrow
        \Spec(A)
    \]
    is a monomorphism of spectral prestacks.

    \item The multiplication morphism
    \[
        B\otimes_A B
        \longrightarrow
        B
    \]
    is an equivalence.
\end{enumerate}
\end{proposition}

\begin{proof}
A morphism $X\to Y$ is a monomorphism if and only if its diagonal is an
equivalence. By \Ref{prop:affine-spectral-fibre-products}, the diagonal of
$\Spec(B)\to\Spec(A)$ is represented by the multiplication map
$B\otimes_A B\to B$.
\end{proof}

\begin{definition}[Affine open immersion]
\label{def:affine-open-immersion-spectral}
A morphism
\[
    \Spec(B)
    \longrightarrow
    \Spec(A)
\]
of affine spectral schemes is an \textbf{affine open immersion} if the
algebra map $A\to B$ is flat, of finite presentation, and satisfies
\[
    B\otimes_A B
    \xrightarrow{\simeq}
    B.
\]
\end{definition}

\begin{proposition}[Flat lifts of classical affine opens]
\label{prop:flat-lift-classical-affine-open}
Let $A\to B$ be a flat morphism of connective commutative ring spectra.
If
\[
    \Spec\pi_0(B)
    \longrightarrow
    \Spec\pi_0(A)
\]
is an open immersion of ordinary affine schemes, then $A\to B$ is of
finite presentation and the multiplication morphism
\[
    B\otimes_A B
    \longrightarrow
    B
\]
is an equivalence.
\end{proposition}

\begin{proof}
This is the affine spectral open-lifting theorem. Flatness identifies the
homotopy groups of $B$ and of all its base changes with extension of the
corresponding $\pi_0(A)$-modules. The classical open immersion is a
finitely presented flat epimorphism; the spectral lifting theorem then
shows that its flat spectral lift is a finitely presented epimorphism.
Equivalently, after a finite principal-open refinement the assertion
reduces to the localization calculation
\[
    \pi_n(A[f^{-1}])
    \simeq
    \pi_n(A)[f^{-1}],
\]
and the local identifications descend over the finite Zariski cover. See
\cite{LurieSAG,KhanBraveNewMotivicI}.
\end{proof}

\begin{theorem}[Classical characterization of affine opens]
\label{thm:classical-characterization-affine-opens}
For a morphism $A\to B$ of connective commutative ring spectra, the
following are equivalent.
\begin{enumerate}[label=(\roman*)]
    \item The induced morphism $\Spec(B)\to\Spec(A)$ is an affine open
    immersion.

    \item The map $A\to B$ is flat and the induced morphism
    \[
        \Spec\pi_0(B)
        \longrightarrow
        \Spec\pi_0(A)
    \]
    is an open immersion of ordinary affine schemes.
\end{enumerate}
\end{theorem}

\begin{proof}
If (i) holds, flat base change identifies the classical truncation of the
diagonal with the classical diagonal. Finite presentation and the
monomorphism condition descend to $\pi_0$, so the classical morphism is a
flat finitely presented monomorphism and hence an open immersion.

Conversely, suppose (ii). Flatness and
\Ref{prop:flat-base-change-classical-truncation} show that
\[
    \pi_0(B\otimes_A B)
    \simeq
    \pi_0(B)
    \otimes_{\pi_0(A)}
    \pi_0(B)
    \simeq
    \pi_0(B).
\]
The flat homotopy-group formula then shows that multiplication induces an
isomorphism on every homotopy group, hence is an equivalence. Finite
presentation, together with the multiplication equivalence, is supplied by
\Ref{prop:flat-lift-classical-affine-open}.
\end{proof}

\begin{proposition}[Permanence of affine open immersions]
\label{prop:affine-open-permanence}
Affine open immersions are stable under equivalence, composition, and base
change.
\end{proposition}

\begin{proof}
Flatness and finite presentation have the required permanence properties
by \Ref{prop:flat-algebra-permanence} and
\Ref{prop:finite-presentation-permanence}. Monomorphisms are stable under
equivalence, composition, and base change. Apply
\Ref{def:affine-open-immersion-spectral}.
\end{proof}

\subsection{Principal opens}

\begin{construction}[Spectral localization]
\label{con:spectral-localization}
Let $A$ be connective and let
\[
    f\in\pi_0(A).
\]
The \textbf{localization of $A$ at $f$} is the initial commutative
$A$-algebra
\[
    A[f^{-1}]
\]
in which the image of $f$ is invertible.
\end{construction}

\begin{proposition}[Homotopy groups of a localization]
\label{prop:homotopy-groups-localization}
For every integer $n$, there is a natural isomorphism
\[
    \pi_n(A[f^{-1}])
    \simeq
    \pi_n(A)[f^{-1}].
\]
In particular, $A\to A[f^{-1}]$ is flat.
\end{proposition}\begin{proof}
The underlying $A$-module of the localization $A[f^{-1}]$ is canonically
equivalent to the filtered colimit
\[
    A
    \xrightarrow{f}
    A
    \xrightarrow{f}
    A
    \xrightarrow{f}
    \cdots
\]
formed in $\Mod_A$, where each arrow is multiplication by $f$ as an
$A$-module map.

Homotopy groups commute with filtered colimits, so for every integer $n$,
\[
\begin{aligned}
    \pi_n(A[f^{-1}])
    &\simeq
    \varinjlim
    \left(
        \pi_n(A)
        \xrightarrow{f}
        \pi_n(A)
        \xrightarrow{f}
        \cdots
    \right)
    \\
    &\simeq
    \pi_n(A)[f^{-1}].
\end{aligned}
\]
In particular,
\[
    \pi_0(A[f^{-1}])
    \simeq
    \pi_0(A)[f^{-1}]
\]
is flat over $\pi_0(A)$, and the displayed homotopy-group formula is
exactly the flatness criterion of
\Ref{thm:spectral-lazard}. Hence
\[
    A\longrightarrow A[f^{-1}]
\]
is flat.
\end{proof}

\begin{proposition}[Principal spectral opens]
\label{prop:principal-spectral-open}
The morphism
\[
    D(f)
    :=
    \Spec(A[f^{-1}])
    \longrightarrow
    \Spec(A)
\]
is an affine open immersion. Its classical truncation is the ordinary
principal open
\[
    D(f)
    \subseteq
    \Spec\pi_0(A).
\]
\end{proposition}
\begin{proof}
The map is flat by
\Ref{prop:homotopy-groups-localization}.

We exhibit its finite presentation. Let
\[
    P
    :=
    \operatorname{Sym}_{\mathbb S}(\mathbb S)
\]
with universal generator \(t\), and give \(A\) its \(P\)-algebra
structure defined by
\[
    t
    \longmapsto
    f.
\]

Work in commutative \(P\)-algebras. Let
\[
    R
    :=
    \operatorname{Sym}_{P}(P)
\]
be the free commutative \(P\)-algebra on one generator \(u\), and let
\[
    Q
    :=
    \operatorname{Sym}_{P}(P)
\]
be the free commutative \(P\)-algebra on one generator \(z\). Form the
pushout
\[
\begin{tikzcd}[column sep=large, row sep=large]
    Q
        \arrow[r, "{z\mapsto tu-1}"]
        \arrow[d, "{z\mapsto0}"']
    &
    R
        \arrow[d]
    \\
    P
        \arrow[r]
    &
    P[t^{-1}].
\end{tikzcd}
\]
The universal property of this pushout identifies its lower-right
corner with the localization of \(P\) in which \(t\) is invertible.
Indeed, a morphism from the pushout to a commutative \(P\)-algebra \(C\)
is equivalently the choice of an element \(u\in\Omega^\infty C\)
together with a homotopy
\[
    tu
    \simeq
    1.
\]
Such data exist exactly when the image of \(t\) is invertible in
\(\pi_0(C)\), and in that case the space of choices is contractible.
Thus the displayed pushout has the universal property of \(P[t^{-1}]\).

The algebra \(R\) is obtained from \(P\) by one free cell attachment.
The morphism
\[
    Q
    \longrightarrow
    P,
    \qquad
    z\longmapsto0,
\]
is the free commutative \(P\)-algebra map induced by the map
\[
    P
    \longrightarrow
    0
\]
of finite free \(P\)-modules, equivalently by extension of scalars from
the map
\[
    \mathbb S
    \longrightarrow
    0
\]
of finite spectra. Hence the displayed pushout attaches one further
finite cell. Therefore \(P[t^{-1}]\) is a finite cell commutative
\(P\)-algebra.

The universal property of localization gives
\[
    A[f^{-1}]
    \simeq
    A\otimes_P P[t^{-1}].
\]
Finite presentation is preserved under base change, so
\[
    A
    \longrightarrow
    A[f^{-1}]
\]
is of finite presentation.

The multiplication morphism
\[
    A[f^{-1}]
    \otimes_A
    A[f^{-1}]
    \longrightarrow
    A[f^{-1}]
\]
is an equivalence by the universal property of localization. Therefore
\[
    D(f)
    =
    \Spec(A[f^{-1}])
    \longrightarrow
    \Spec(A)
\]
is an affine open immersion.

Finally,
\Ref{prop:homotopy-groups-localization}
gives
\[
    \pi_0(A[f^{-1}])
    \simeq
    \pi_0(A)[f^{-1}],
\]
so the classical truncation is the ordinary principal open
\[
    D(f)
    \subseteq
    \Spec\pi_0(A).
\]
\end{proof}

\begin{proposition}[Principal opens form a basis]
\label{prop:principal-opens-basis}
Every affine open immersion into $\Spec(A)$ is covered by principal open
subschemes $D(f)$ with $f\in\pi_0(A)$.
\end{proposition}
\begin{proof}
Let
\[
    j
    \colon
    U=\Spec(B)
    \longrightarrow
    X=\Spec(A)
\]
be an affine open immersion. By
\Ref{thm:classical-characterization-affine-opens},
the classical truncation
\[
    j_{\mathrm{cl}}
    \colon
    U_{\mathrm{cl}}
    =
    \Spec\pi_0(B)
    \longrightarrow
    X_{\mathrm{cl}}
    =
    \Spec\pi_0(A)
\]
is an open immersion of ordinary affine schemes.

Ordinary principal opens form a basis. Hence, for every point of
\(U_{\mathrm{cl}}\), there is an element
\[
    f\in\pi_0(A)
\]
such that
\[
    D(f)
    \subseteq
    X_{\mathrm{cl}}
\]
contains the point and is contained in \(U_{\mathrm{cl}}\).

Form the spectral fibre product
\[
    P
    :=
    U
    \times_X
    \Spec(A[f^{-1}]).
\]
By
\Ref{prop:affine-spectral-fibre-products},
it is affine:
\[
    P
    \simeq
    \Spec
    \bigl(
        B\otimes_AA[f^{-1}]
    \bigr).
\]
The projection
\[
    p
    \colon
    P
    \longrightarrow
    \Spec(A[f^{-1}])
\]
is the base change of the affine open immersion \(U\to X\). It is
therefore an affine open immersion and, in particular, flat.

Classical truncation preserves this fibre product by
\Ref{prop:flat-base-change-classical-truncation}. Since the ordinary
principal open \(D(f)\) is contained in \(U_{\mathrm{cl}}\), one obtains
\[
\begin{aligned}
    P_{\mathrm{cl}}
    &\simeq
    U_{\mathrm{cl}}
    \times_{X_{\mathrm{cl}}}
    D(f)
    \\
    &\simeq
    D(f).
\end{aligned}
\]
Consequently, the algebra morphism representing \(p\),
\[
    A[f^{-1}]
    \longrightarrow
    B\otimes_AA[f^{-1}],
\]
induces an isomorphism on degree-zero homotopy groups.

Because this algebra morphism is flat, the flat homotopy-group formula
of \Ref{prop:homotopy-groups-flat-algebra} gives, for every integer \(n\),
\[
\begin{aligned}
    \pi_n
    \bigl(
        B\otimes_AA[f^{-1}]
    \bigr)
    &\simeq
    \pi_n
    \bigl(
        A[f^{-1}]
    \bigr)
    \otimes_{\pi_0(A[f^{-1}])}
    \pi_0
    \bigl(
        B\otimes_AA[f^{-1}]
    \bigr)
    \\
    &\simeq
    \pi_n
    \bigl(
        A[f^{-1}]
    \bigr).
\end{aligned}
\]
Thus the representing algebra morphism is an equivalence, and therefore
\[
    P
    \xrightarrow{\simeq}
    \Spec(A[f^{-1}]).
\]

The inverse equivalence identifies the principal spectral open
\[
    D(f)
    =
    \Spec(A[f^{-1}])
\]
with an open spectral subscheme of \(U\). Choosing such an \(f\) around
every point of \(U_{\mathrm{cl}}\) gives a family of principal spectral
opens whose classical truncations cover \(U_{\mathrm{cl}}\). Hence, in
the terminology of
\Ref{def:affine-zariski-cover},
this family is an affine spectral Zariski covering family of \(U\).
\end{proof}

\subsection{The affine spectral Zariski topology}

\begin{definition}[Affine Zariski covering family]
\label{def:affine-zariski-cover}
A small family of affine open immersions
\[
    \{\Spec(B_i)\to\Spec(A)\}_{i\in I}
\]
is an \textbf{affine spectral Zariski covering family} if the ordinary
opens
\[
    \{\Spec\pi_0(B_i)\to\Spec\pi_0(A)\}_{i\in I}
\]
cover the topological space $\Spec\pi_0(A)$.
\end{definition}

\begin{proposition}[The spectral Zariski pretopology]
\label{prop:spectral-zariski-pretopology}
Affine spectral Zariski covering families define a Grothendieck
pretopology on $\Aff_{\Sp}$.
\end{proposition}

\begin{proof}
The identity family covers. Base change stability follows from
\Ref{prop:affine-open-permanence} and the base-change compatibility of
classical truncation in
\Ref{prop:flat-base-change-classical-truncation}. Transitivity reduces to
transitivity of ordinary Zariski covers on classical truncations.
\end{proof}

\begin{definition}[Spectral Zariski topology]
\label{def:spectral-zariski-topology}
The \textbf{spectral Zariski topology}
\[
    \tau_{\mathrm{Zar}}
\]
on $\Aff_{\Sp}$ is the Grothendieck topology generated by the
pretopology of \Ref{prop:spectral-zariski-pretopology}.
\end{definition}

\begin{proposition}[Joint conservativity criterion]
\label{prop:zariski-cover-joint-conservativity}
A family of affine open immersions
\[
    \{\Spec(B_i)\to\Spec(A)\}_{i\in I}
\]
is a Zariski covering family if and only if the pullback functors
\[
    \{-\otimes_A B_i\}_{i\in I}
\]
are jointly conservative on $\Mod_A$.
\end{proposition}

\begin{proof}
Suppose first that the family covers. If $M\otimes_AB_i\simeq0$ for every
$i$, then
\[
    \pi_n(M)
    \otimes_{\pi_0(A)}
    \pi_0(B_i)
    \simeq
    0
\]
for every $n$, by flatness. The ordinary quasi-coherent module
$\pi_n(M)$ therefore vanishes after restriction to a Zariski cover of
$\Spec\pi_0(A)$ and hence vanishes. Thus all homotopy groups of $M$ vanish,
so $M\simeq0$.

Conversely, if the classical opens do not cover, choose a point in the
closed complement and a nonzero residue-field module supported at that
point. Its Eilenberg--Mac Lane lift is an $A$-module annihilated by every
$-\otimes_AB_i$. Joint conservativity therefore fails.
\end{proof}

\subsection{Descendability and Zariski descent}

\begin{definition}[Thick tensor ideal]
\label{def:spectral-thick-tensor-ideal}
Let $\mathcal C^\otimes$ be a stable symmetric monoidal
$\infty$-category. A full subcategory
\[
    \mathcal I
    \subseteq
    \mathcal C
\]
is a \textbf{thick tensor ideal} if it is closed under finite limits,
finite colimits, retracts, and tensoring with arbitrary objects of
$\mathcal C$.
\end{definition}

\begin{definition}[Descendable morphism]
\label{def:spectral-descendable-morphism}
A morphism
\[
    A
    \longrightarrow
    B
\]
of commutative ring spectra is \textbf{descendable} if $B$ generates
$\Mod_A$ as a thick tensor ideal.
\end{definition}

\begin{proposition}[Permanence of descendability]
\label{prop:spectral-descendability-permanence}
Descendable morphisms are stable under equivalence, base change, and
composition. If $A\to B\to C$ is a factorization and $A\to C$ is
descendable, then $A\to B$ is descendable.
\end{proposition}

\begin{proof}
Base change follows because extension of scalars is exact and strong
symmetric monoidal. The remaining assertions are the permanence results of
\cite[Proposition~3.24]{MathewGaloisStableHomotopy}.
\end{proof}

\begin{lemma}[Modules over finite product algebras]
\label{lem:spectral-modules-finite-products}
For a finite family of commutative ring spectra $\{B_i\}_{i\in I}$, there
is a canonical symmetric monoidal equivalence
\[
    \Mod_{\prod_iB_i}
    \simeq
    \prod_i\Mod_{B_i}.
\]
\end{lemma}

\begin{proof}
The finite product of commutative ring spectra has underlying finite
biproduct of spectra. Its orthogonal idempotents split every module into
its components. Extension of scalars along the projections and finite
direct sum of restrictions of scalars define mutually inverse symmetric
monoidal functors, exactly as for ordinary product rings.
\end{proof}

\begin{proposition}[Finite Zariski covers are descendable]
\label{prop:finite-zariski-covers-descendable}
Let
\[
    \{\Spec(B_i)\to\Spec(A)\}_{i\in I}
\]
be a finite affine spectral Zariski covering family and put
\[
    B
    :=
    \prod_{i\in I}B_i.
\]
Then
\[
    A
    \longrightarrow
    B
\]
is descendable.
\end{proposition}

\begin{proof}
Choose a finite principal refinement
\[
    \{D(f_j)\to\Spec(A)\}_{j=1}^r
\]
subordinate to the given cover, where
\[
    f_1,\ldots,f_r
    \in
    \pi_0(A)
\]
generate the unit ideal. Put
\[
    A_j
    :=
    A[f_j^{-1}]
\]
and
\[
    C
    :=
    \prod_{j=1}^rA_j.
\]
The refinement determines a factorization
\[
    A
    \longrightarrow
    B
    \longrightarrow
    C.
\]
By the intermediate-algebra property of
\Ref{prop:spectral-descendability-permanence}, it is enough to prove that
$A\to C$ is descendable.

Consider the finite augmented alternating \v{C}ech complex of $A$-modules
\[
\begin{aligned}
    A
    &\longrightarrow
    \bigoplus_jA_j
    \longrightarrow
    \bigoplus_{j_0<j_1}
    A_{j_0}\otimes_AA_{j_1}
    \longrightarrow
    \cdots
    \\
    &\longrightarrow
    A_1\otimes_A\cdots\otimes_AA_r.
\end{aligned}
\]
For every integer $n$, applying $\pi_n$ identifies this complex with the
ordinary augmented \v{C}ech complex
\[
\begin{aligned}
    \pi_n(A)
    &\longrightarrow
    \bigoplus_j\pi_n(A)[f_j^{-1}]
    \longrightarrow
    \bigoplus_{j_0<j_1}
    \pi_n(A)[(f_{j_0}f_{j_1})^{-1}]
    \longrightarrow
    \cdots.
\end{aligned}
\]
This ordinary complex is exact because the principal opens
$D(f_j)$ cover $\Spec\pi_0(A)$. Hence the spectral augmented complex is
exact.

Every nonaugmented term is a finite direct sum of tensor products of the
modules $A_j$. The finite direct sum
\[
    \bigoplus_jA_j
\]
is the underlying $A$-module of $C$. Consequently every nonaugmented term
belongs to the thick tensor ideal generated by $C$.

The augmented complex is finite and reconstructs $A$ by iterated fibres.
A thick tensor ideal is closed under finite fibres. Therefore $A$ belongs
to the thick tensor ideal generated by $C$, and $A\to C$ is descendable.

Since
\[
    A
    \longrightarrow
    B
    \longrightarrow
    C
\]
is a factorization and $A\to C$ is descendable,
\Ref{prop:spectral-descendability-permanence} implies that $A\to B$ is
descendable.
\end{proof}

\begin{theorem}[Amitsur convergence]
\label{thm:spectral-amitsur-convergence}
If $A\to B$ is descendable, then the augmented Amitsur object
\[
    A
    \longrightarrow
    B
    \triplearrows
    B\otimes_A B
    \cdots
\]
is a limit diagram in $\CAlg(\Sp)$. In particular,
\[
    A
    \simeq
    \operatorname*{Tot}_{[n]\in\Delta}
    B^{\otimes_A(n+1)}.
\]
\end{theorem}

\begin{proof}
Regard $B$ as a commutative algebra object of $\Mod_A$. Descendability
means that $B$ generates the tensor unit. The descendability theorem
identifies the augmented cobar object as a limit diagram in $\Mod_A$; see
\cite[Proposition~3.20]{MathewGaloisStableHomotopy}. The forgetful functor
from commutative $A$-algebras to $A$-modules creates limits, so the same
object is a limit in commutative $A$-algebras and hence in
$\CAlg(\Sp)$.
\end{proof}

\begin{theorem}[Affine module descent]
\label{thm:spectral-affine-module-descent}
If $A\to B$ is descendable, then extension of scalars induces an
equivalence in $\CAlg(\Cat_\infty)$
\[
    \Mod_A
    \xrightarrow{\simeq}
    \operatorname*{Tot}_{[n]\in\Delta}
    \Mod_{B^{\otimes_A(n+1)}}.
\]
\end{theorem}

\begin{proof}
The extension--restriction adjunction is comonadic for a descendable map,
and its cobar resolution is the displayed Amitsur diagram; see
\cite[Proposition~3.22]{MathewGaloisStableHomotopy}. Every transition
functor is strong symmetric monoidal, so the comparison is an equivalence
of symmetric monoidal $\infty$-categories.
\end{proof}

\begin{theorem}[Affine algebra descent]
\label{thm:spectral-affine-algebra-descent}
If $A\to B$ is descendable, then base change induces an equivalence
\[
    \CAlg(\Sp)_{A/}
    \xrightarrow{\simeq}
    \operatorname*{Tot}_{[n]\in\Delta}
    \CAlg(\Sp)_{B^{\otimes_A(n+1)}/}.
\]
\end{theorem}

\begin{proof}
By \Ref{thm:spectral-affine-module-descent}, extension of scalars gives a
symmetric monoidal equivalence
\[
    \Mod_A
    \xrightarrow{\simeq}
    \operatorname*{Tot}_{[n]\in\Delta}
    \Mod_{B^{\otimes_A(n+1)}}.
\]
The functor assigning to a symmetric monoidal $\infty$-category its
$\infty$-category of commutative algebra objects preserves limits.
Applying it to the displayed equivalence gives
\[
    \CAlg(\Mod_A)
    \xrightarrow{\simeq}
    \operatorname*{Tot}_{[n]\in\Delta}
    \CAlg
    \bigl(
        \Mod_{B^{\otimes_A(n+1)}}
    \bigr).
\]
Using
\Ref{prop:relative-commutative-spectral-algebras} in every cosimplicial
degree identifies this equivalence with
\[
    \CAlg(\Sp)_{A/}
    \xrightarrow{\simeq}
    \operatorname*{Tot}_{[n]\in\Delta}
    \CAlg(\Sp)_{B^{\otimes_A(n+1)}/}.
\]
\end{proof}

\begin{corollary}[Connective algebra descent for finite Zariski covers]
\label{cor:connective-zariski-algebra-descent}
Let
\[
    \{\Spec(B_i)\to\Spec(A)\}_{i\in I}
\]
be a finite affine spectral Zariski covering family and put
\[
    B
    :=
    \prod_{i\in I}B_i.
\]
Then the equivalence of
\Ref{thm:spectral-affine-algebra-descent} restricts to an equivalence
\[
    \CAlg(\Sp)^{\mathrm{cn}}_{A/}
    \xrightarrow{\simeq}
    \operatorname*{Tot}_{[n]\in\Delta}
    \CAlg(\Sp)^{\mathrm{cn}}_
    {B^{\otimes_A(n+1)}/}.
\]
\end{corollary}

\begin{proof}
The morphism $A\to B$ is descendable by
\Ref{prop:finite-zariski-covers-descendable}. It is also flat: each
$B_i$ is flat over $A$, and the finite product $B$ is the finite direct
sum of the $B_i$ as an $A$-module.

Every Amitsur term
\[
    B^{\otimes_A(n+1)}
\]
is connective. Base change along $A\to B$ is $t$-exact and conservative
on modules.

Let a descent datum in the right side consist of connective commutative
algebras. By the unrestricted algebra-descent theorem, it descends to a
commutative $A$-algebra $C$. Its base change
\[
    C\otimes_AB
\]
is connective. Since $-\otimes_AB$ is $t$-exact,
\[
    \tau_{\leq -1}(C)\otimes_AB
    \simeq
    \tau_{\leq -1}(C\otimes_AB)
    \simeq
    0.
\]
Conservativity of $-\otimes_AB$ implies
\[
    \tau_{\leq -1}(C)
    \simeq
    0.
\]
Thus $C$ is connective. Full faithfulness is inherited from the
unrestricted descent equivalence.
\end{proof}

\begin{theorem}[Subcanonicity of the spectral Zariski topology]
\label{thm:spectral-zariski-subcanonical}
Every representable spectral prestack is a sheaf for
$\tau_{\mathrm{Zar}}$.
\end{theorem}

\begin{proof}
It is enough to verify \v{C}ech descent on finite affine covers, since every
affine Zariski cover admits a finite principal refinement locally on a
quasi-compact affine. Let $\{\Spec(B_i)\to\Spec(A)\}$ be a finite cover and
put $B=\prod_iB_i$. By
\Ref{prop:finite-zariski-covers-descendable}, $A\to B$ is descendable.
For every connective commutative ring spectrum $C$,
\Ref{thm:spectral-amitsur-convergence} gives
\[
    \Map_{\CAlg(\Sp)}(C,A)
    \simeq
    \operatorname*{Tot}_{[n]\in\Delta}
    \Map_{\CAlg(\Sp)}
    \bigl(
        C,
        B^{\otimes_A(n+1)}
    \bigr).
\]
The finite-product decomposition of the Amitsur terms identifies this
with the componentwise \v{C}ech descent condition for the covering family.
Hence the representable $h_C$ is a sheaf.
\end{proof}

\begin{theorem}[Zariski descent for affine quasi-coherent modules]
\label{thm:affine-qcoh-zariski-descent}
Let
\[
    U:=\coprod_i\Spec(B_i)
    \longrightarrow
    \Spec(A)
\]
be a finite affine spectral Zariski cover, and let $U_\bullet$ be its
\v{C}ech nerve. Then the comparison functor is an equivalence
\[
    \Mod_A
    \xrightarrow{\simeq}
    \operatorname*{Tot}_{[n]\in\Delta}
    \QCoh(U_n).
\]
\end{theorem}

\begin{proof}
Put $B=\prod_iB_i$. The map $A\to B$ is descendable by
\Ref{prop:finite-zariski-covers-descendable}. Apply
\Ref{thm:spectral-affine-module-descent}. The decomposition of modules over
finite product algebras from
\Ref{lem:spectral-modules-finite-products} identifies the Amitsur diagram
with the componentwise \v{C}ech diagram of the family.
\end{proof}

\section{Spectral stacks and spectral schemes}
\label{sec:spectral-stacks-schemes}

\begin{definition}[Spectral Zariski stack]
\label{def:spectral-zariski-stack}
A \textbf{spectral Zariski stack} is a sheaf of spaces on
$\Aff_{\Sp}$ for the spectral Zariski topology. The
$\infty$-category of spectral Zariski stacks is denoted
\[
    \operatorname{Stk}^{\mathrm{sp}}_{\mathrm{Zar}}
    :=
    \Sh_{\tau_{\mathrm{Zar}}}
    (\Aff_{\Sp}).
\]
\end{definition}

By \Ref{thm:spectral-zariski-subcanonical}, the Yoneda embedding factors as
\[
\begin{tikzcd}[column sep=large]
    \Aff_{\Sp}
        \arrow[r, hook, "h"]
    &
    \operatorname{Stk}^{\mathrm{sp}}_{\mathrm{Zar}}
        \arrow[r, hook]
    &
    \PSh(\Aff_{\Sp}).
\end{tikzcd}
\]

\subsection{Representable open immersions}

\begin{definition}[Open immersion of spectral stacks]
\label{def:open-immersion-spectral-stacks}
Let
\[
    j
    \colon
    U
    \longrightarrow
    X
\]
be a morphism of spectral Zariski stacks.

Suppose first that $X$ is an affine spectral scheme. Then $j$ is an
\textbf{open immersion} if:
\begin{enumerate}[label=(\roman*)]
    \item $j$ is a monomorphism;

    \item there exists a small family of affine spectral schemes
    $\{U_\alpha\}_{\alpha\in I}$ and morphisms
    \[
        U_\alpha
        \longrightarrow
        U
    \]
    such that
    \[
        \coprod_{\alpha\in I}U_\alpha
        \longrightarrow
        U
    \]
    is an effective epimorphism and every composite
    \[
        U_\alpha
        \longrightarrow
        U
        \longrightarrow
        X
    \]
    is an affine open immersion in the sense of
    \Ref{def:affine-open-immersion-spectral}.
\end{enumerate}

For a general spectral Zariski stack $X$, the morphism $j$ is an open
immersion if, for every affine spectral scheme $T$ and every morphism
$T\to X$, the pullback
\[
    U\times_XT
    \longrightarrow
    T
\]
is an open immersion in the preceding sense.
\end{definition}

\begin{proposition}[Permanence of open immersions]
\label{prop:global-open-immersion-permanence}
Open immersions of spectral Zariski stacks are stable under equivalence,
composition, and base change.
\end{proposition}

\begin{proof}
Stability under equivalence is immediate. Stability under base change is
built into the definition.

Consider composable open immersions
\[
    U
    \xrightarrow{j}
    V
    \xrightarrow{k}
    X.
\]
It is enough to prove that their composite is open after base change
along an affine spectral scheme $T\to X$. Thus assume that $X=T$ is
affine.

Choose an affine open atlas
\[
    \coprod_\alpha V_\alpha
    \longrightarrow
    V
\]
such that each composite
\[
    V_\alpha
    \longrightarrow
    V
    \longrightarrow
    T
\]
is an affine open immersion. For each $\alpha$, the pullback
\[
    U_\alpha
    :=
    U\times_VV_\alpha
\]
is open in the affine spectral scheme $V_\alpha$. Choose an affine open
atlas
\[
    \coprod_\beta U_{\alpha\beta}
    \longrightarrow
    U_\alpha
\]
whose composites
\[
    U_{\alpha\beta}
    \longrightarrow
    V_\alpha
\]
are affine open immersions.

By composition stability of affine open immersions,
\[
    U_{\alpha\beta}
    \longrightarrow
    T
\]
is an affine open immersion. The family
\[
    \coprod_{\alpha,\beta}
    U_{\alpha\beta}
    \longrightarrow
    U
\]
is an effective epimorphism, because effective epimorphisms are stable
under base change and composition in an $\infty$-topos. The composite
$U\to T$ is also a monomorphism. Hence it is an open immersion.
\end{proof}

\subsection{Spectral schemes}

\begin{definition}[Spectral scheme]
\label{def:spectral-scheme}
A \textbf{spectral scheme} is a spectral Zariski stack $X$ for which there
exists a small family of affine spectral schemes $\{U_i\}_{i\in I}$ and a
morphism
\[
    \coprod_{i\in I}U_i
    \longrightarrow
    X
\]
such that:
\begin{enumerate}[label=(\roman*)]
    \item the displayed morphism is an effective epimorphism;

    \item every component $U_i\to X$ is an open immersion.
\end{enumerate}
Such a family is an \textbf{affine spectral Zariski atlas} of $X$.
\end{definition}

The full subcategory of
\[
    \operatorname{Stk}^{\mathrm{sp}}_{\mathrm{Zar}}
\]
spanned by spectral schemes is denoted
\[
    \operatorname{Sch}^{\mathrm{sp}}.
\]

\begin{proposition}[Affine spectral schemes are spectral schemes]
\label{prop:affines-are-spectral-schemes}
Every affine spectral scheme is a spectral scheme.
\end{proposition}

\begin{proof}
Use the singleton atlas given by the identity map.
\end{proof}

\begin{proposition}[Coproducts of spectral schemes]
\label{prop:spectral-schemes-coproducts}
Small coproducts of spectral schemes exist and are spectral schemes.
\end{proposition}

\begin{proof}
Coproducts exist in the ambient $\infty$-topos of spectral Zariski stacks.
Taking the coproducts of affine atlases gives an effective epimorphism by
a coproduct of affine open immersions.
\end{proof}

\begin{theorem}[Fibre products of spectral schemes]
\label{thm:spectral-schemes-fibre-products}
Let
\[
    X
    \longrightarrow
    Z
    \longleftarrow
    Y
\]
be morphisms of spectral schemes. Then the fibre product
\[
    X\times_ZY
\]
is a spectral scheme.
\end{theorem}

\begin{proof}
Choose affine Zariski atlases $U\to X$, $V\to Y$, and $T\to Z$. Pulling
$T\to Z$ back to each component of $U$ and $V$ produces affine open
covers because open immersions are representable and stable under base
change. Refining by these covers reduces the fibre product to fibre
products of affine spectral schemes over affines, which are affine by
\Ref{prop:affine-spectral-fibre-products}. The resulting family maps by
open immersions to $X\times_ZY$ and is an effective epimorphism because
effective epimorphisms are universal in an $\infty$-topos.
\end{proof}

\begin{corollary}[Finite limits of spectral schemes]
\label{cor:spectral-schemes-finite-limits}
The $\infty$-category of spectral schemes admits finite limits, computed
in spectral Zariski stacks.
\end{corollary}

\begin{proof}
The terminal object is $\Spec(\mathbb S)$, and binary fibre products exist
by \Ref{thm:spectral-schemes-fibre-products}.
\end{proof}

\subsection{Classical truncation of spectral schemes}

\begin{theorem}[Global classical truncation]
\label{thm:global-classical-truncation}
The affine classical truncation functor glues to a functor
\[
    (-)_{\mathrm{cl}}
    \colon
    \operatorname{Sch}^{\mathrm{sp}}
    \longrightarrow
    \operatorname{Sch}
\]
from spectral schemes to ordinary schemes. For every spectral scheme $X$,
there is a natural morphism
\[
    X_{\mathrm{cl}}
    \longrightarrow
    X
\]
whose restriction to every affine open $\Spec(A)\subseteq X$ is
\[
    \Spec\pi_0(A)
    \longrightarrow
    \Spec(A).
\]
\end{theorem}

\begin{proof}
Choose an affine Zariski atlas $U\to X$. Its \v{C}ech nerve consists of
affine spectral schemes after refinement by affine opens. Classical
truncation preserves the open base changes occurring in this diagram by
\Ref{prop:flat-base-change-classical-truncation}. The classical affine
truncations therefore form a Zariski groupoid in ordinary affine schemes.
Its colimit is an ordinary scheme $X_{\mathrm{cl}}$. The natural affine
maps to $X$ are compatible on intersections and descend to the displayed
morphism. Independence of the atlas follows from common refinement.
\end{proof}

\begin{lemma}[Principal classical opens determine spectral principal opens]
\label{lem:principal-classical-open-determines-spectral-open}
Let \(A\) be a connective commutative ring spectrum, and let
\[
    j
    \colon
    U
    \longrightarrow
    \Spec(A)
\]
be an open immersion of spectral schemes. Suppose that, as an open
subscheme of \(\Spec\pi_0(A)\), the classical truncation of \(U\) is
\[
    U_{\mathrm{cl}}
    =
    D(f)
\]
for some \(f\in\pi_0(A)\). Then there is a canonical equivalence over
\(\Spec(A)\)
\[
    U
    \simeq
    \Spec(A[f^{-1}]).
\]
\end{lemma}

\begin{proof}
Put
\[
    X
    :=
    \Spec(A)
\]
and
\[
    V
    :=
    \Spec(A[f^{-1}]).
\]
By
\Ref{prop:principal-spectral-open},
the morphism
\[
    V
    \longrightarrow
    X
\]
is an open immersion, and
\[
    V_{\mathrm{cl}}
    =
    D(f)
    =
    U_{\mathrm{cl}}
\]
as open subschemes of \(X_{\mathrm{cl}}\).

We first record a general observation. Let
\[
    W
    \longrightarrow
    T
\]
be an open immersion of spectral schemes which induces an equivalence
\[
    W_{\mathrm{cl}}
    \xrightarrow{\simeq}
    T_{\mathrm{cl}}.
\]
Then \(W\to T\) is an equivalence.

Indeed, this assertion is local on \(T\). After base change along an
affine open
\[
    T_0
    \longrightarrow
    T,
\]
we may suppose that \(T=T_0\) is affine. By the definition of an open
immersion into an affine spectral scheme, there is a family of affine
spectral schemes \(\{W_\alpha\}_{\alpha\in I}\) and an effective
epimorphism
\[
    \coprod_{\alpha\in I}W_\alpha
    \longrightarrow
    W
\]
such that every composite
\[
    W_\alpha
    \longrightarrow
    W
    \longrightarrow
    T
\]
is an affine open immersion. Since
\[
    W_{\mathrm{cl}}
    \xrightarrow{\simeq}
    T_{\mathrm{cl}},
\]
the corresponding classical affine opens cover \(T_{\mathrm{cl}}\).
Hence
\[
    \{W_\alpha\to T\}_{\alpha\in I}
\]
is an affine spectral Zariski covering family. It follows that
\[
    \coprod_{\alpha\in I}W_\alpha
    \longrightarrow
    T
\]
is an effective epimorphism.

The morphism \(W\to T\) is a monomorphism, and the preceding effective
epimorphism factors through it. Therefore \(W\to T\) is itself an
effective epimorphism. A morphism in an \(\infty\)-topos which is both a
monomorphism and an effective epimorphism is an equivalence. Thus
\[
    W
    \xrightarrow{\simeq}
    T.
\]

Now form the fibre product
\[
    P
    :=
    U\times_XV.
\]
Both projections
\[
    P
    \longrightarrow
    U
    \qquad\text{and}\qquad
    P
    \longrightarrow
    V
\]
are open immersions. By
\Ref{prop:global-classical-truncation-fibre-products},
\[
\begin{aligned}
    P_{\mathrm{cl}}
    &\simeq
    U_{\mathrm{cl}}
    \times_{X_{\mathrm{cl}}}
    V_{\mathrm{cl}}
    \\
    &\simeq
    D(f)
    \times_{X_{\mathrm{cl}}}
    D(f)
    \\
    &\simeq
    D(f).
\end{aligned}
\]
Consequently, each projection induces an equivalence on classical
truncations. By the observation above, both projections are
equivalences. Hence
\[
    U
    \simeq
    P
    \simeq
    V
    =
    \Spec(A[f^{-1}])
\]
over \(X=\Spec(A)\).
\end{proof}

\begin{proposition}[Classical truncation and fibre products]
\label{prop:global-classical-truncation-fibre-products}
For every diagram of spectral schemes
\[
    X
    \longrightarrow
    Z
    \longleftarrow
    Y,
\]
there is a canonical equivalence of ordinary schemes
\[
    (X\times_ZY)_{\mathrm{cl}}
    \simeq
    X_{\mathrm{cl}}
    \times_{Z_{\mathrm{cl}}}
    Y_{\mathrm{cl}}.
\]
\end{proposition}

\begin{proof}
The assertion is local on all three schemes for the Zariski topology.
After choosing compatible affine open covers, it reduces to a span
\[
    \Spec(B)
    \longrightarrow
    \Spec(A)
    \longleftarrow
    \Spec(C).
\]
The spectral fibre product is $\Spec(B\otimes_A C)$, and
\Ref{prop:flat-base-change-classical-truncation} gives
\[
    \pi_0(B\otimes_A C)
    \simeq
    \pi_0(B)
    \otimes_{\pi_0(A)}
    \pi_0(C).
\]
The resulting affine equivalences agree on overlaps and therefore glue.
\end{proof}

\begin{theorem}[Affineness is detected classically]
\label{thm:affineness-detected-classically}
A spectral scheme $X$ is affine if and only if its classical truncation
$X_{\mathrm{cl}}$ is affine.
\end{theorem}
\begin{proof}
If \(X\) is affine, then \(X_{\mathrm{cl}}\) is affine by definition.

Conversely, suppose
\[
    X_{\mathrm{cl}}
    \simeq
    \Spec(R).
\]
Choose an affine spectral Zariski atlas
\[
    \{W_\alpha\to X\}_{\alpha\in J}.
\]
The ordinary affine opens
\[
    (W_\alpha)_{\mathrm{cl}}
\]
cover \(\Spec(R)\). Since ordinary principal opens form a basis and
\(\Spec(R)\) is quasi-compact, there are elements
\[
    f_1,\ldots,f_r
    \in
    R
\]
which generate the unit ideal and such that every principal open
\[
    D(f_i)
\]
is contained in some
\[
    (W_{\alpha(i)})_{\mathrm{cl}}.
\]

Inside \(W_{\alpha(i)}\), take the principal spectral open determined by
the restriction of \(f_i\). Denote its source by
\[
    U_i
    =
    \Spec(B_i).
\]
Its composite
\[
    U_i
    \longrightarrow
    W_{\alpha(i)}
    \longrightarrow
    X
\]
is an open immersion, and its classical truncation is the principal open
\[
    (U_i)_{\mathrm{cl}}
    =
    D(f_i)
    \subseteq
    \Spec(R).
\]
Since the \(f_i\) generate the unit ideal, the \(U_i\) form a finite
affine open cover of \(X\).

For a tuple
\[
    \mathbf i
    =
    (i_0,\ldots,i_p),
\]
put
\[
    U_{\mathbf i}
    :=
    U_{i_0}
    \times_X
    \cdots
    \times_X
    U_{i_p}.
\]
The morphism
\[
    U_{\mathbf i}
    \longrightarrow
    U_{i_0}
\]
is an open immersion. Its classical truncation is
\[
\begin{aligned}
    (U_{\mathbf i})_{\mathrm{cl}}
    &\simeq
    D(f_{i_0})
    \times_{\Spec(R)}
    \cdots
    \times_{\Spec(R)}
    D(f_{i_p})
    \\
    &=
    D(f_{\mathbf i}),
\end{aligned}
\]
where
\[
    f_{\mathbf i}
    :=
    f_{i_0}\cdots f_{i_p}.
\]
By
\Ref{lem:principal-classical-open-determines-spectral-open},
the open spectral subscheme \(U_{\mathbf i}\) is the corresponding
principal spectral open of \(U_{i_0}\). In particular, it is affine.
Write
\[
    U_{\mathbf i}
    =
    \Spec(B_{\mathbf i}).
\]

For every integer \(n\geq0\), let \(B_n^\bullet\) be the cosimplicial
commutative ring spectrum whose degree-\(p\) term is
\[
    B_n^p
    :=
    \prod_{\mathbf i\in\{1,\ldots,r\}^{p+1}}
    \tau_{\leq n}B_{\mathbf i},
\]
with coface and codegeneracy maps induced by the \v{C}ech nerve of the
cover. Define
\[
    A_n
    :=
    \operatorname*{Tot}_{[p]\in\Delta}
    B_n^p
\]
in \(\CAlg(\Sp)\).

For every integer \(q\geq0\), the modules \(\pi_q(B_i)\), together with
their localization identifications on the intersections, define an
ordinary quasi-coherent module
\[
    \mathcal M_q
\]
on \(\Spec(R)\). Put
\[
    M_q
    :=
    \Gamma
    \bigl(
        \Spec(R),
        \mathcal M_q
    \bigr).
\]
Flatness of the open restrictions and the homotopy-group formula for
localization give
\[
    \pi_q(B_{\mathbf i})
    \simeq
    M_q[f_{\mathbf i}^{-1}].
\]

The totalization spectral sequence for the underlying cosimplicial
spectra has
\[
    E_1^{p,q}
    =
    \pi_q(B_n^p)
    \Longrightarrow
    \pi_{q-p}(A_n).
\]
It has only the rows
\[
    0\leq q\leq n,
\]
and therefore converges strongly. For fixed \(q\), its \(E_1\)-complex
is the ordinary augmented \v{C}ech complex of the quasi-coherent module
\(\mathcal M_q\) for the finite principal cover
\[
    \{D(f_i)\}_{i=1}^r.
\]
Since \(\Spec(R)\) is affine, that augmented complex is exact. Hence
\[
    E_2^{p,q}
    \simeq
    \begin{cases}
        M_q,&p=0,\\
        0,&p>0.
    \end{cases}
\]
It follows that
\[
    \pi_q(A_n)
    \simeq
    M_q
    \qquad
    (0\leq q\leq n),
\]
and all other homotopy groups of \(A_n\) vanish. In particular,
\[
    \pi_0(A_n)
    \simeq
    R.
\]

The totalization cone gives, for each \(i\), a morphism
\[
    A_n
    \longrightarrow
    \tau_{\leq n}B_i.
\]
The image of \(f_i\in\pi_0(A_n)\simeq R\) is invertible in the target.
Consequently, this morphism factors through a map
\[
    A_n[f_i^{-1}]
    \longrightarrow
    \tau_{\leq n}B_i.
\]
On homotopy groups this is the canonical isomorphism
\[
    M_q[f_i^{-1}]
    \xrightarrow{\simeq}
    \pi_q(B_i)
    \qquad
    (0\leq q\leq n).
\]
Both sides are \(n\)-truncated and connective, so
\[
    A_n[f_i^{-1}]
    \xrightarrow{\simeq}
    \tau_{\leq n}B_i.
\]

The truncation morphisms induce a tower
\[
    \cdots
    \longrightarrow
    A_n
    \longrightarrow
    A_{n-1}
    \longrightarrow
    \cdots
    \longrightarrow
    A_0.
\]
Define
\[
    A
    :=
    \varprojlim_nA_n
\]
in \(\CAlg(\Sp)\).

For every fixed \(q\geq0\), the inverse system
\[
    \{\pi_q(A_n)\}_n
\]
is eventually constant with value \(M_q\). The Milnor exact sequence
therefore gives
\[
    \pi_q(A)
    \simeq
    M_q
\]
for every \(q\geq0\). For \(q<0\), both the inverse limit and the relevant
derived inverse limit vanish, again because the homotopy groups of the
tower are eventually constant. Thus
\[
    \pi_q(A)=0
    \qquad
    (q<0),
\]
so \(A\) is connective.

The compatible maps
\[
    A_n
    \longrightarrow
    \tau_{\leq n}B_i,
\]
together with Postnikov completeness of \(B_i\), induce a morphism
\[
    A
    \longrightarrow
    B_i.
\]
The element \(f_i\) is invertible in \(B_i\), so this morphism factors
through
\[
    A[f_i^{-1}]
    \longrightarrow
    B_i.
\]
On every homotopy group it induces the canonical isomorphism
\[
    M_q[f_i^{-1}]
    \xrightarrow{\simeq}
    \pi_q(B_i).
\]
Hence
\[
    A[f_i^{-1}]
    \xrightarrow{\simeq}
    B_i.
\]

The maps
\[
    \Spec(B_i)
    \longrightarrow
    \Spec(A)
\]
agree on all iterated intersections because they arise from the same
totalization cones. They therefore glue to a morphism
\[
    X
    \longrightarrow
    \Spec(A).
\]
On the finite affine open cover \(\{U_i\}\), this morphism identifies
\(U_i\) with the principal open
\[
    D(f_i)
    \subseteq
    \Spec(A).
\]
Since
\[
    \pi_0(A)
    \simeq
    R
\]
and the elements \(f_1,\ldots,f_r\) generate the unit ideal, these
principal opens cover \(\Spec(A)\). The glued morphism is therefore an
equivalence. Thus
\[
    X
    \simeq
    \Spec(A).
\]
See \cite{LurieSAG}.
\end{proof}

\begin{definition}[Underlying topological space]
\label{def:underlying-space-spectral-scheme}
For a spectral scheme $X$, define
\[
    |X|
    :=
    |X_{\mathrm{cl}}|.
\]
\end{definition}

\begin{proposition}[Open subsets and classical truncation]
\label{prop:opens-classical-truncation}
The assignment
\[
    U
    \longmapsto
    U_{\mathrm{cl}}
\]
induces an equivalence between the poset of open spectral subschemes of
$X$ and the poset of open subschemes of $X_{\mathrm{cl}}$.
\end{proposition}

\begin{proof}
The assertion is local on $X$. For $X=\Spec(A)$ it follows from
\Ref{thm:classical-characterization-affine-opens} and the fact that
principal opens form a basis. The affine identifications glue uniquely.
\end{proof}

\subsection{Quasi-compactness and quasi-separatedness}

\begin{definition}[Quasi-compact spectral scheme]
\label{def:quasi-compact-spectral-scheme}
A spectral scheme $X$ is \textbf{quasi-compact} if it admits a finite
affine open cover. A morphism $f:X\to Y$ is quasi-compact if the pullback
over every affine open of $Y$ is quasi-compact.
\end{definition}

\begin{definition}[Quasi-separated morphisms and spectral schemes]
\label{def:quasi-separated-spectral-scheme}
A morphism
\[
    f
    \colon
    X
    \longrightarrow
    Y
\]
of spectral schemes is \textbf{quasi-separated} if its diagonal
\[
    \Delta_f
    \colon
    X
    \longrightarrow
    X\times_YX
\]
is quasi-compact.

A spectral scheme $X$ is \textbf{quasi-separated} if its structural
morphism
\[
    X
    \longrightarrow
    \Spec(\mathbb S)
\]
is quasi-separated. Equivalently, the absolute diagonal
\[
    X
    \longrightarrow
    X\times X
\]
is quasi-compact.

A spectral scheme is \textbf{qcqs} if it is quasi-compact and
quasi-separated.
\end{definition}

\begin{proposition}[Classical detection of qcqs properties]
\label{prop:classical-detection-qcqs}
A spectral scheme is quasi-compact, quasi-separated, or qcqs if and only
if its classical truncation has the corresponding property.
\end{proposition}

\begin{proof}
Affine open covers and their intersections correspond under
\Ref{prop:opens-classical-truncation}. Quasi-compactness and
quasi-separatedness are therefore detected by the same finite-cover
conditions on $X$ and $X_{\mathrm{cl}}$.
\end{proof}

\section{Quasi-coherent modules on spectral schemes and relative spectra}
\label{sec:qcoh-spectral-schemes-relative-spec}

The formal quasi-coherent category of
\Ref{def:qcoh-spectral-prestack} acquires its geometric interpretation on
spectral schemes through Zariski descent.

\subsection{Descent and the global \texorpdfstring{$t$}{t}-structure}

\begin{theorem}[Atlas computation of quasi-coherent modules]
\label{thm:qcoh-atlas-computation}
Let $X$ be a spectral scheme and let
\[
    U
    \longrightarrow
    X
\]
be an affine Zariski atlas. Write $U_\bullet$ for its \v{C}ech nerve in
spectral schemes. Then the canonical comparison functor is a symmetric
monoidal equivalence
\[
    \QCoh(X)
    \xrightarrow{\simeq}
    \operatorname*{Tot}_{[n]\in\Delta}
    \QCoh(U_n).
\]
\end{theorem}

\begin{proof}
An object of the totalization on the right is a quasi-coherent module on
$U$, together with coherent descent equivalences on
\[
    U\times_XU,
    \quad
    U\times_XU\times_XU,
    \quad\ldots.
\]
Restriction of a quasi-coherent module on $X$ gives such data, producing
the comparison functor.

We construct an inverse. Let
\[
    v
    \colon
    V=\Spec(A)
    \longrightarrow
    X
\]
be an arbitrary affine point. The pullback
\[
    U_V
    :=
    U\times_XV
    \longrightarrow
    V
\]
is a Zariski covering morphism. Since $V$ is quasi-compact, this cover
admits a finite affine refinement
\[
    \coprod_{i=1}^r
    V_i
    \longrightarrow
    V,
    \qquad
    V_i=\Spec(A_i).
\]
The given descent datum on $U_\bullet$ restricts to a descent datum on the
\v{C}ech nerve of this finite affine cover. By
\Ref{thm:affine-qcoh-zariski-descent}, it descends uniquely to an
$A$-module
\[
    M_V
    \in
    \Mod_A.
\]

The construction is independent of the chosen finite affine refinement.
Indeed, two refinements admit a common finite affine refinement, and
full faithfulness in affine module descent identifies the two descended
modules by a canonical equivalence. The space of such identifications is
contractible.

Now let
\[
    V'=\Spec(A')
    \longrightarrow
    V=\Spec(A)
\]
be any morphism of affine points over $X$. Base-changing a finite affine
refinement of $U_V\to V$ gives a finite affine refinement of
$U_{V'}\to V'$. The modules
\[
    M_V\otimes_AA'
    \qquad\text{and}\qquad
    M_{V'}
\]
descend the same pulled-back datum. Affine module descent therefore gives
a canonical equivalence
\[
    M_V\otimes_AA'
    \simeq
    M_{V'}.
\]
Uniqueness in descent supplies the cocycle and all higher compatibility
conditions.

Thus the modules $M_V$, as $V$ ranges over all affine points of $X$,
define an object
\[
    M
    \in
    \varprojlim_
    {(\Spec(A)\to X)
    \in
    (\Aff_{\Sp}/X)^{\op}}
    \Mod_A
    =
    \QCoh(X).
\]
Affine module descent shows that restricting $M$ back to $U_\bullet$
recovers the original descent datum. Conversely, starting with an object
of $\QCoh(X)$ and applying the two constructions recovers its restrictions
to every affine point and hence the original object.

All constructions use strong symmetric monoidal pullback functors, and
all limits are formed in
\[
    \CAlg(\Pr^L_{\mathrm{st}}).
\]
The resulting equivalence is therefore symmetric monoidal.
\end{proof}

\begin{corollary}[Zariski locality of quasi-coherent modules]
\label{cor:qcoh-zariski-locality}
A quasi-coherent module on a spectral scheme is equivalently a family of
modules on the members of an affine Zariski atlas, together with coherent
descent equivalences on all iterated intersections.
\end{corollary}

\begin{proof}
Unfold the totalization in \Ref{thm:qcoh-atlas-computation}.
\end{proof}

\begin{lemma}[Computation on the affine-open basis]
\label{lem:qcoh-affine-open-basis}
Let $X$ be a spectral scheme, and let
\[
    \operatorname{AffOp}(X)
\]
denote the $\infty$-category of affine open spectral subschemes of $X$.
Restriction induces a symmetric monoidal equivalence
\[
    \QCoh(X)
    \xrightarrow{\simeq}
    \varprojlim_
    {V\in\operatorname{AffOp}(X)^{\op}}
    \QCoh(V).
\]
\end{lemma}

\begin{proof}
Restriction defines the displayed functor. Conversely, a compatible
family of modules on the affine opens of $X$ restricts, on any affine
Zariski atlas and all its iterated intersections, to coherent descent
data. The atlas equivalence of
\Ref{thm:qcoh-atlas-computation} glues these data to an object of
$\QCoh(X)$. The two constructions are inverse because affine-open
restrictions jointly detect equivalences. Symmetric monoidality follows
from symmetric monoidality of restriction and from formation of the limit
in $\CAlg(\Pr^L_{\mathrm{st}})$.
\end{proof}

\begin{definition}[Flat morphisms and flat quasi-coherent modules]
\label{def:flat-spectral-scheme-morphism}
A morphism
\[
    f
    \colon
    Y
    \longrightarrow
    X
\]
of spectral schemes is \textbf{flat} if, for every affine spectral scheme
\[
    U
    =
    \Spec(A)
    \longrightarrow
    X
\]
and every affine open
\[
    V
    =
    \Spec(B)
    \subseteq
    Y\times_XU,
\]
the induced morphism $A\to B$ is flat.

A quasi-coherent module $\mathcal F\in\QCoh(X)$ is \textbf{flat} if, for
every affine open $U=\Spec(A)\subseteq X$, the corresponding $A$-module
$\mathcal F|_U$ is flat.
\end{definition}

\begin{proposition}[Zariski locality of flatness]
\label{prop:global-flatness-zariski-local}
Flatness of morphisms of spectral schemes is local on the source and the
target for the spectral Zariski topology. Consequently, in
\Ref{def:flat-spectral-scheme-morphism} it is enough to verify flatness on
an affine Zariski atlas of the target and an affine Zariski atlas of each
pullback to the source.
\end{proposition}

\begin{proof}
It is enough to prove locality on the source in the affine case. Let
$A\to B$ be a morphism of connective commutative ring spectra and let
\[
    \{\Spec(B_i)\to\Spec(B)\}_{i\in I}
\]
be a Zariski covering family such that every composite $A\to B_i$ is
flat. Let $N\in\Mod_{A,\leq0}$ and put
\[
    M
    :=
    N\otimes_AB.
\]
For every $i$,
\[
    M\otimes_BB_i
    \simeq
    N\otimes_AB_i
    \in
    \Mod_{B_i,\leq0}.
\]
The maps $B\to B_i$ are flat, so their extension-of-scalars functors are
$t$-exact. Hence
\[
    \tau_{\geq1}(M)\otimes_BB_i
    \simeq
    \tau_{\geq1}
    \bigl(
        M\otimes_BB_i
    \bigr)
    \simeq
    0.
\]
The functors $-\otimes_BB_i$ are jointly conservative by
\Ref{prop:zariski-cover-joint-conservativity}. Therefore
\[
    \tau_{\geq1}(M)
    \simeq
    0,
\]
so $M$ is coconnective. Extension of scalars along $A\to B$ is already
right $t$-exact, and is therefore $t$-exact. Thus $A\to B$ is flat.

Locality on the target follows by base change, locality on the source,
and refinement by affine open covers. The final assertion follows.
\end{proof}

\begin{proposition}[Zariski locality of flat modules]
\label{prop:flat-module-zariski-local}
Let \(A\) be a connective commutative ring spectrum, let
\[
    \{\Spec(A_i)\to\Spec(A)\}_{i\in I}
\]
be an affine spectral Zariski covering family, and let
\[
    M
    \in
    \Mod_A.
\]
Then \(M\) is flat over \(A\) if and only if
\[
    M\otimes_AA_i
\]
is flat over \(A_i\) for every \(i\).
\end{proposition}

\begin{proof}
If \(M\) is flat over \(A\), then
\[
    M\otimes_AA_i
\]
is flat over \(A_i\) by base-change permanence of flat modules from
\Ref{prop:flat-module-permanence}.

Conversely, suppose that
\[
    M_i
    :=
    M\otimes_AA_i
\]
is flat over \(A_i\) for every \(i\). Each map
\[
    A
    \longrightarrow
    A_i
\]
is flat because it represents an affine open immersion. Hence extension
of scalars
\[
    -\otimes_AA_i
    \colon
    \Mod_A
    \longrightarrow
    \Mod_{A_i}
\]
is \(t\)-exact.

We first prove that \(M\) is connective. Since every \(M_i\) is flat, it
is connective. Therefore
\[
\begin{aligned}
    \tau_{\leq-1}(M)\otimes_AA_i
    &\simeq
    \tau_{\leq-1}
    \bigl(
        M\otimes_AA_i
    \bigr)
    \\
    &\simeq
    \tau_{\leq-1}(M_i)
    \\
    &\simeq
    0.
\end{aligned}
\]
The extension-of-scalars functors
\[
    \{-\otimes_AA_i\}_{i\in I}
\]
are jointly conservative by
\Ref{prop:zariski-cover-joint-conservativity}.
Hence
\[
    \tau_{\leq-1}(M)
    \simeq
    0,
\]
so \(M\) is connective.

It remains to prove that tensoring with \(M\) preserves coconnective
modules. Let
\[
    N
    \in
    \Mod_{A,\leq0}.
\]
For every \(i\), there is a natural equivalence
\[
\begin{aligned}
    (N\otimes_AM)\otimes_AA_i
    &\simeq
    \bigl(
        N\otimes_AA_i
    \bigr)
    \otimes_{A_i}
    \bigl(
        M\otimes_AA_i
    \bigr)
    \\
    &\simeq
    \bigl(
        N\otimes_AA_i
    \bigr)
    \otimes_{A_i}
    M_i.
\end{aligned}
\]
Since \(A\to A_i\) is flat,
\[
    N\otimes_AA_i
    \in
    \Mod_{A_i,\leq0}.
\]
Since \(M_i\) is flat over \(A_i\), it follows that
\[
    (N\otimes_AM)\otimes_AA_i
    \in
    \Mod_{A_i,\leq0}.
\]
Using \(t\)-exactness of extension of scalars once more gives
\[
\begin{aligned}
    \tau_{\geq1}(N\otimes_AM)\otimes_AA_i
    &\simeq
    \tau_{\geq1}
    \bigl(
        (N\otimes_AM)\otimes_AA_i
    \bigr)
    \\
    &\simeq
    0.
\end{aligned}
\]
Joint conservativity therefore implies
\[
    \tau_{\geq1}(N\otimes_AM)
    \simeq
    0.
\]
Thus
\[
    N\otimes_AM
    \in
    \Mod_{A,\leq0}.
\]

Because \(M\) is connective, tensoring with \(M\) is already right
\(t\)-exact. The preceding argument proves left \(t\)-exactness.
Therefore
\[
    -\otimes_AM
\]
is \(t\)-exact, and \(M\) is flat over \(A\).
\end{proof}

\begin{theorem}[Canonical $t$-structure on global quasi-coherent modules]
\label{thm:qcoh-spectral-scheme-t-structure}
Let $X$ be a spectral scheme. The full subcategories
\[
    \QCoh(X)_{\geq0}
\]
and
\[
    \QCoh(X)_{\leq0}
\]
of quasi-coherent modules which are respectively connective and
coconnective on every affine open define an accessible, left complete,
and right complete $t$-structure on $\QCoh(X)$.

Pullback along every morphism of spectral schemes is right $t$-exact.
Pullback along a flat morphism is $t$-exact.
\end{theorem}

\begin{proof}
By \Ref{lem:qcoh-affine-open-basis},
\[
    \QCoh(X)
    \simeq
    \varprojlim_
    {V\in\operatorname{AffOp}(X)^{\op}}
    \Mod_{\Gamma(V,\mathcal O_V)}.
\]
Every transition functor is restriction along an open immersion. Open
immersions are flat, so the transition functors are $t$-exact.

For a compatible family
\[
    \mathcal F
    =
    \{\mathcal F_V\}_V
\]
in the displayed limit, apply the affine truncation functors
componentwise. Since every transition functor is $t$-exact, the objects
$\tau_{\geq0}\mathcal F_V$ and $\tau_{\leq-1}\mathcal F_V$ remain
compatible and define objects of $\QCoh(X)$. Their componentwise
truncation triangles assemble to a fibre and cofibre sequence
\[
    \tau_{\geq0}\mathcal F
    \longrightarrow
    \mathcal F
    \longrightarrow
    \tau_{\leq-1}\mathcal F.
\]
Orthogonality follows because mapping spectra in the limit are limits of
the affine mapping spectra. Thus these subcategories define the asserted
$t$-structure.

To verify accessibility, choose a regular cardinal $\kappa$ such that
all affine module categories in the chosen small affine-open diagram, all
restriction functors between them, and all affine truncation functors are
$\kappa$-accessible. The componentwise truncation functors preserve the
limiting compatibility conditions because every restriction functor is
$t$-exact. They therefore induce $\kappa$-accessible truncation functors
on $\QCoh(X)$. Thus the global $t$-structure is accessible.

For left completeness, the affine-open-basis equivalence and the
componentwise definition of the $t$-structure give, for every integer
$n$, an equivalence
\[
    \QCoh(X)_{\leq n}
    \simeq
    \varprojlim_
    {V\in\operatorname{AffOp}(X)^{\op}}
    \Mod_{\Gamma(V,\mathcal O_V),\leq n}.
\]
Every affine module category is left complete. Since limits commute with
limits in $\Cat_\infty$, there are natural equivalences
\[
\begin{aligned}
    \QCoh(X)
    &\simeq
    \varprojlim_V
    \Mod_{\Gamma(V,\mathcal O_V)}
    \\
    &\simeq
    \varprojlim_V
    \varprojlim_n
    \Mod_{\Gamma(V,\mathcal O_V),\leq n}
    \\
    &\simeq
    \varprojlim_n
    \varprojlim_V
    \Mod_{\Gamma(V,\mathcal O_V),\leq n}
    \\
    &\simeq
    \varprojlim_n
    \QCoh(X)_{\leq n}.
\end{aligned}
\]
Under these identifications, the resulting functor is the canonical
Postnikov-completion functor. Thus the global $t$-structure is left
complete.

For right completeness, consider the canonical morphism
\[
    \varinjlim_n
    \tau_{\geq-n}\mathcal F
    \longrightarrow
    \mathcal F.
\]
Restriction to every affine open preserves the filtered colimit and
identifies this morphism with the right-completeness equivalence in the
corresponding affine module category. Since affine-open restrictions
jointly detect equivalences, the displayed morphism is an equivalence.

Let $f:Y\to X$ be any morphism. On affine charts, $f^*$ is extension of
scalars along a morphism of connective commutative ring spectra and is
right $t$-exact by
\Ref{prop:base-change-exact-right-t-exact}. Therefore the global
pullback is right $t$-exact.

If $f$ is flat, then its affine base changes are flat. On affine charts,
extension of scalars is $t$-exact by the definition of flatness. Hence the
global pullback is $t$-exact.
\end{proof}

\begin{theorem}[Heart and classical truncation]
\label{thm:qcoh-heart-classical-truncation}
For every spectral scheme $X$, restriction of scalars along
\[
    i_X
    \colon
    X_{\mathrm{cl}}
    \longrightarrow
    X
\]
induces an equivalence of abelian categories
\[
    i_{X,*}
    \colon
    \QCoh(X_{\mathrm{cl}})^\heartsuit
    \xrightarrow{\simeq}
    \QCoh(X)^\heartsuit.
\]
Its inverse is
\[
    \mathcal F
    \longmapsto
    \pi_0\bigl(i_X^*\mathcal F\bigr).
\]
Equivalently,
\[
    \QCoh(X)^\heartsuit
    \simeq
    \operatorname{QCoh}_{\mathrm{ab}}(X_{\mathrm{cl}}),
\]
where the right side denotes the ordinary abelian category of
quasi-coherent modules on $X_{\mathrm{cl}}$.
\end{theorem}

\begin{proof}
The assertion is local for the Zariski topology. Let
\[
    X=\Spec(A).
\]
The map $i_X$ is induced by the canonical algebra morphism
\[
    A
    \longrightarrow
    H\pi_0(A).
\]
Restriction of scalars gives
\[
    \Mod_{H\pi_0(A)}^\heartsuit
    \longrightarrow
    \Mod_A^\heartsuit.
\]
Every discrete $A$-module is determined by its degree-zero homotopy group,
and its action factors uniquely through $H\pi_0(A)$. Hence restriction of
scalars is fully faithful and essentially surjective.

For a discrete $A$-module $M$, the degree-zero edge morphism gives
\[
    \pi_0
    \left(
        H\pi_0(A)\otimes_AM
    \right)
    \simeq
    \pi_0(M).
\]
Thus the inverse equivalence on the heart is
\[
    M
    \longmapsto
    \pi_0
    \left(
        H\pi_0(A)\otimes_AM
    \right).
\]
These affine equivalences are natural under restriction to open
subschemes and therefore glue through
\Ref{lem:qcoh-affine-open-basis}.

The use of $\pi_0$ is essential: the raw pullback $i_X^*$ need not carry a
heart object to a heart object.
\end{proof}

\subsection{Quasi-coherent commutative algebras}

\begin{definition}[Quasi-coherent commutative algebra]
\label{def:quasi-coherent-commutative-algebra}
For a spectral scheme \(X\), define
\[
    \operatorname{QCAlg}(X)
    :=
    \CAlg\bigl(\QCoh(X)\bigr).
\]
Let
\[
    U_X
    \colon
    \operatorname{QCAlg}(X)
    \longrightarrow
    \QCoh(X)
\]
denote the forgetful functor. The full subcategory of algebra objects
whose underlying quasi-coherent modules are connective is denoted
\[
    \operatorname{QCAlg}(X)^{\mathrm{cn}}
    :=
    \left\{
        \mathcal A
        \in
        \operatorname{QCAlg}(X)
        \,\middle|\,
        U_X(\mathcal A)
        \in
        \QCoh(X)_{\geq0}
    \right\}.
\]
No \(t\)-structure on the generally nonstable
\(\infty\)-category \(\operatorname{QCAlg}(X)\) is being asserted.
\end{definition}

\begin{proposition}[Pullback of quasi-coherent algebras]
\label{prop:pullback-qcalg}
Every morphism $f:X\to Y$ of spectral schemes induces a functor
\[
    f^*
    \colon
    \operatorname{QCAlg}(Y)
    \longrightarrow
    \operatorname{QCAlg}(X).
\]
It preserves connective algebra objects, and for composable morphisms
there are canonical coherent equivalences
\[
    (g\circ f)^*
    \simeq
    f^*g^*.
\]
\end{proposition}

\begin{proof}
The quasi-coherent pullback functor is strong symmetric monoidal by
\Ref{prop:qcoh-pullback}, hence carries commutative algebra objects to
commutative algebra objects. Right $t$-exactness gives preservation of
connectivity. Coherence is inherited from the functoriality of
quasi-coherent pullback.
\end{proof}

\subsection{Relative spectra}

\begin{theorem}[Relative spectrum construction]
\label{thm:relative-spectrum-construction}
Let $X$ be a spectral scheme. There is a functor
\[
    \Spec_X
    \colon
    \operatorname{QCAlg}(X)^{\mathrm{cn},\op}
    \longrightarrow
    \operatorname{Sch}^{\mathrm{sp}}_{/X}
\]
characterized by the following properties.
\begin{enumerate}[label=(\roman*)]
    \item If $X=\Spec(A)$ and $\mathcal A$ corresponds to a connective
    commutative $A$-algebra $B$, then
    \[
        \Spec_X(\mathcal A)
        \simeq
        \Spec(B).
    \]

    \item For every morphism $f:Y\to X$, there is a natural equivalence
    \[
        \Map_X
        \bigl(
            Y,
            \Spec_X(\mathcal A)
        \bigr)
        \simeq
        \Map_{\operatorname{QCAlg}(Y)}
        \bigl(
            f^*\mathcal A,
            \mathcal O_Y
        \bigr).
    \]
\end{enumerate}
\end{theorem}

\begin{proof}
Let
\[
    \mathcal A
    \in
    \operatorname{QCAlg}(X)^{\mathrm{cn}}.
\]
For every affine spectral scheme $T$, define
\[
    F_{\mathcal A}(T)
    :=
    \operatorname*{colim}_{f\in X(T)}
    \Map_{\operatorname{QCAlg}(T)}
    \bigl(
        f^*\mathcal A,
        \mathcal O_T
    \bigr),
\]
where the colimit is taken over the $\infty$-groupoid $X(T)$.
Equivalently, $F_{\mathcal A}(T)$ is the total space of the
space-valued family over $X(T)$ whose fibre over a morphism
\[
    f
    \colon
    T
    \longrightarrow
    X
\]
is
\[
    \Map_{\operatorname{QCAlg}(T)}
    \bigl(
        f^*\mathcal A,
        \mathcal O_T
    \bigr).
\]

For a morphism of affine spectral schemes
\[
    g
    \colon
    T'
    \longrightarrow
    T,
\]
pullback carries a point
\[
    \alpha
    \colon
    f^*\mathcal A
    \longrightarrow
    \mathcal O_T
\]
to the point
\[
    g^*\alpha
    \colon
    (f\circ g)^*\mathcal A
    \longrightarrow
    \mathcal O_{T'}.
\]
These maps define a spectral prestack
\[
    F_{\mathcal A}
    \longrightarrow
    X.
\]
A morphism
\[
    \mathcal A
    \longrightarrow
    \mathcal A'
\]
induces, by precomposition, a morphism
\[
    F_{\mathcal A'}
    \longrightarrow
    F_{\mathcal A}.
\]
Thus the construction is contravariantly functorial in $\mathcal A$.

We first prove that $F_{\mathcal A}$ is a spectral Zariski stack. Let
\[
    \{T_i\to T\}_{i\in I}
\]
be a finite affine spectral Zariski covering family, with $T$ and every
$T_i$ affine. For a tuple
\[
    \mathbf i
    =
    (i_0,\ldots,i_n),
\]
write
\[
    T_{\mathbf i}
    :=
    T_{i_0}
    \times_T
    \cdots
    \times_T
    T_{i_n}.
\]
Every $T_{\mathbf i}$ is affine by
\Ref{prop:affine-spectral-fibre-products}.

Since $X$ is a spectral Zariski stack, there is a natural equivalence
\[
    X(T)
    \simeq
    \operatorname*{Tot}_{[n]\in\Delta}
    \prod_{\mathbf i\in I^{n+1}}
    X(T_{\mathbf i}).
\]
By affine quasi-coherent descent and preservation of limits by the
commutative-algebra construction, there is an equivalence
\[
    \operatorname{QCAlg}(T)
    \simeq
    \operatorname*{Tot}_{[n]\in\Delta}
    \prod_{\mathbf i\in I^{n+1}}
    \operatorname{QCAlg}(T_{\mathbf i}).
\]

There is a canonical comparison morphism
\[
    F_{\mathcal A}(T)
    \longrightarrow
    \operatorname*{Tot}_{[n]\in\Delta}
    \prod_{\mathbf i\in I^{n+1}}
    F_{\mathcal A}(T_{\mathbf i})
\]
lying over the preceding equivalence for $X(T)$.

Fix a point
\[
    f
    \in
    X(T),
\]
and let
\[
    f_{\mathbf i}
    \colon
    T_{\mathbf i}
    \longrightarrow
    X
\]
denote its restrictions. The fibre of the comparison morphism over the
compatible family $\{f_{\mathbf i}\}$ is
\[
\begin{aligned}
    &
    \Map_{\operatorname{QCAlg}(T)}
    \bigl(
        f^*\mathcal A,
        \mathcal O_T
    \bigr)
    \\
    &\qquad\longrightarrow
    \operatorname*{Tot}_{[n]\in\Delta}
    \prod_{\mathbf i\in I^{n+1}}
    \Map_{\operatorname{QCAlg}(T_{\mathbf i})}
    \bigl(
        f_{\mathbf i}^*\mathcal A,
        \mathcal O_{T_{\mathbf i}}
    \bigr).
\end{aligned}
\]
Under the equivalence
\[
    \operatorname{QCAlg}(T)
    \simeq
    \operatorname*{Tot}_{[n]\in\Delta}
    \prod_{\mathbf i\in I^{n+1}}
    \operatorname{QCAlg}(T_{\mathbf i}),
\]
the objects $f^*\mathcal A$ and $\mathcal O_T$ correspond respectively
to the compatible families
\[
    \{f_{\mathbf i}^*\mathcal A\}_{\mathbf i}
    \qquad\text{and}\qquad
    \{\mathcal O_{T_{\mathbf i}}\}_{\mathbf i}.
\]
Mapping spaces in a limit of $\infty$-categories are the corresponding
limits of mapping spaces. Therefore the displayed morphism on fibres is
an equivalence.

The comparison morphism is consequently an equivalence on its base and
on every fibre, and hence is an equivalence. Thus $F_{\mathcal A}$
satisfies \v{C}ech descent for finite affine spectral Zariski covering
families.

Every affine spectral Zariski covering family admits a finite principal
refinement, and finite affine covers generate the spectral Zariski
topology. It follows that $F_{\mathcal A}$ is a spectral Zariski stack.

Let
\[
    U
    =
    \Spec(R)
    \longrightarrow
    X
\]
be an affine point. Under
\Ref{prop:connective-relative-algebras},
the restriction
\[
    \mathcal A|_U
\]
corresponds to a connective commutative $R$-algebra $B$.

For every affine spectral scheme
\[
    T
    =
    \Spec(C)
    \longrightarrow
    U,
\]
there are natural equivalences
\[
\begin{aligned}
    \Map_{\operatorname{QCAlg}(T)}
    \bigl(
        (\mathcal A|_U)|_T,
        \mathcal O_T
    \bigr)
    &\simeq
    \Map_{\CAlg_C}
    \bigl(
        B\otimes_R C,
        C
    \bigr)
    \\
    &\simeq
    \Map_{\CAlg_R}(B,C)
    \\
    &\simeq
    \Map_U
    \bigl(
        T,
        \Spec(B)
    \bigr).
\end{aligned}
\]
These equivalences are natural in $T$ and identify the fibre product
over $U$:
\[
    F_{\mathcal A}
    \times_X
    U
    \simeq
    \Spec(B).
\]

Choose an affine spectral Zariski atlas
\[
    \{U_i\to X\}_{i\in I}.
\]
Then
\[
    \coprod_{i\in I}
    \bigl(
        F_{\mathcal A}
        \times_X
        U_i
    \bigr)
    \longrightarrow
    F_{\mathcal A}
\]
is the base change of the effective epimorphism
\[
    \coprod_{i\in I}U_i
    \longrightarrow
    X.
\]
It is therefore an effective epimorphism. Each morphism
\[
    F_{\mathcal A}
    \times_X
    U_i
    \longrightarrow
    F_{\mathcal A}
\]
is an open immersion by base change, and every source is affine by the
preceding calculation. Hence $F_{\mathcal A}$ is a spectral scheme over
$X$.

Define
\[
    \Spec_X(\mathcal A)
    :=
    F_{\mathcal A}.
\]
The affine calculation proves property~\textup{(i)}.

It remains to prove property~\textup{(ii)} for an arbitrary spectral
scheme over $X$. On the Zariski site of spectral schemes over $X$, define
presheaves of spaces
\[
    \Phi
    \bigl(
        f:Y\to X
    \bigr)
    :=
    \Map_X
    \bigl(
        Y,
        \Spec_X(\mathcal A)
    \bigr)
\]
and
\[
    \Psi
    \bigl(
        f:Y\to X
    \bigr)
    :=
    \Map_{\operatorname{QCAlg}(Y)}
    \bigl(
        f^*\mathcal A,
        \mathcal O_Y
    \bigr).
\]

The presheaf $\Phi$ is a Zariski sheaf because
$\Spec_X(\mathcal A)$ is a spectral Zariski stack. Explicitly, if
\[
    Y_0
    \longrightarrow
    Y
\]
is an affine Zariski atlas with \v{C}ech nerve $Y_\bullet$, then
\[
    \Phi(Y)
    \simeq
    \operatorname*{Tot}_{[n]\in\Delta}
    \Phi(Y_n).
\]

The presheaf $\Psi$ is also a Zariski sheaf. Indeed,
\Ref{thm:qcoh-atlas-computation}
and preservation of limits by $\CAlg$ give
\[
    \operatorname{QCAlg}(Y)
    \simeq
    \operatorname*{Tot}_{[n]\in\Delta}
    \operatorname{QCAlg}(Y_n).
\]
Under this equivalence, $f^*\mathcal A$ and $\mathcal O_Y$ correspond to
their restrictions to the \v{C}ech nerve. Since mapping spaces in a limit
of $\infty$-categories are computed as limits, one obtains
\[
    \Psi(Y)
    \simeq
    \operatorname*{Tot}_{[n]\in\Delta}
    \Psi(Y_n).
\]

For affine $Y$, the defining construction of $F_{\mathcal A}$ gives a
natural equivalence
\[
    \Phi(Y)
    \simeq
    \Psi(Y).
\]
Affine spectral schemes over $X$ form a basis for the Zariski topology
on spectral schemes over $X$. A natural equivalence between Zariski
sheaves on this affine basis therefore extends uniquely to a natural
equivalence on the entire site. Consequently, for every spectral scheme
$f:Y\to X$,
\[
    \Map_X
    \bigl(
        Y,
        \Spec_X(\mathcal A)
    \bigr)
    \simeq
    \Map_{\operatorname{QCAlg}(Y)}
    \bigl(
        f^*\mathcal A,
        \mathcal O_Y
    \bigr).
\]
This proves property~\textup{(ii)}.

Contravariant functoriality in $\mathcal A$ was constructed above, and
uniqueness follows from the universal property. See
\cite{LurieSAG,KhanBraveNewMotivicI}.
\end{proof}

\begin{proposition}[Base change for relative spectra]
\label{prop:relative-spectrum-base-change}
Let $f:Y\to X$ be a morphism of spectral schemes and let
$\mathcal A\in\operatorname{QCAlg}(X)^{\mathrm{cn}}$. There is a canonical
equivalence over $Y$
\[
    \Spec_X(\mathcal A)
    \times_X
    Y
    \simeq
    \Spec_Y(f^*\mathcal A).
\]
\end{proposition}

\begin{proof}
For every $T\to Y$, the universal property of the left side gives
\[
\begin{aligned}
    \Map_Y
    \bigl(
        T,
        \Spec_X(\mathcal A)\times_XY
    \bigr)
    &\simeq
    \Map_X
    \bigl(
        T,
        \Spec_X(\mathcal A)
    \bigr)
    \\
    &\simeq
    \Map_{\operatorname{QCAlg}(T)}
    \bigl(
        \mathcal A|_T,
        \mathcal O_T
    \bigr)
    \\
    &\simeq
    \Map_Y
    \bigl(
        T,
        \Spec_Y(f^*\mathcal A)
    \bigr).
\end{aligned}
\]
Yoneda gives the equivalence.
\end{proof}

\begin{definition}[Affine morphism of spectral schemes]
\label{def:affine-morphism-spectral-schemes}
A morphism $p:Y\to X$ of spectral schemes is \textbf{affine} if, for every
affine open $U\to X$, the pullback
\[
    Y\times_XU
\]
is an affine spectral scheme.
\end{definition}

\begin{theorem}[Classification of affine morphisms]
\label{thm:classification-affine-morphisms}
For every spectral scheme $X$, the relative spectrum construction induces
an equivalence
\[
    \operatorname{QCAlg}(X)^{\mathrm{cn},\op}
    \xrightarrow{\simeq}
    \bigl(
        \operatorname{Sch}^{\mathrm{sp}}_{/X}
    \bigr)_{\mathrm{aff}}
\]
onto the full subcategory of affine morphisms over $X$.
\end{theorem}

\begin{proof}
The assertion is affine-local on $X$. For $X=\Spec(A)$ it is exactly
\Ref{prop:connective-relative-algebras}. The affine equivalences are
compatible on overlaps by base change and therefore glue through Zariski
descent.
\end{proof}

\section{Closed immersions, complements, and closed gluing}
\label{sec:closed-immersions-complements}

\subsection{Closed immersions}

\begin{definition}[Affine closed immersion]
\label{def:affine-closed-immersion-spectral}
A morphism
\[
    \Spec(B)
    \longrightarrow
    \Spec(A)
\]
of affine spectral schemes is a \textbf{closed immersion} if the induced
ordinary ring homomorphism
\[
    \pi_0(A)
    \longrightarrow
    \pi_0(B)
\]
is surjective.
\end{definition}

\begin{definition}[Closed immersion of spectral schemes]
\label{def:closed-immersion-spectral-schemes}
A morphism
\[
    i
    \colon
    Z
    \longrightarrow
    X
\]
of spectral schemes is a \textbf{closed immersion} if, after base change
along every affine open $U\to X$, the pullback
\[
    Z\times_XU
    \longrightarrow
    U
\]
is an affine closed immersion.
\end{definition}

\begin{proposition}[Classical detection of closed immersions]
\label{prop:closed-immersions-classical-detection}
A morphism $i:Z\to X$ of spectral schemes is a closed immersion if and
only if
\[
    i_{\mathrm{cl}}
    \colon
    Z_{\mathrm{cl}}
    \longrightarrow
    X_{\mathrm{cl}}
\]
is a closed immersion of ordinary schemes.
\end{proposition}

\begin{proof}
If $i$ is a closed immersion, the assertion follows directly on affine
opens from the definition.

Conversely, suppose $i_{\mathrm{cl}}$ is a closed immersion. The assertion
is local on $X$, so let
\[
    X=\Spec(A).
\]
Then $Z_{\mathrm{cl}}$ is a closed subscheme of the affine scheme
$\Spec\pi_0(A)$ and is therefore affine. By
\Ref{thm:affineness-detected-classically}, the spectral scheme $Z$ is
affine, say
\[
    Z=\Spec(B).
\]
The ordinary ring homomorphism
\[
    \pi_0(A)
    \longrightarrow
    \pi_0(B)
\]
is surjective because it represents the given classical closed
immersion. Hence
\[
    \Spec(B)
    \longrightarrow
    \Spec(A)
\]
is an affine closed immersion by
\Ref{def:affine-closed-immersion-spectral}. The local statements glue,
so $i$ is a closed immersion.
\end{proof}

\begin{proposition}[Permanence of closed immersions]
\label{prop:closed-immersion-permanence}
Closed immersions of spectral schemes are stable under equivalence,
composition, and base change.
\end{proposition}

\begin{proof}
The assertion is local on the target and is detected on classical
truncations by
\Ref{prop:closed-immersions-classical-detection}. Apply the corresponding
permanence properties for ordinary closed immersions.
\end{proof}

\begin{proposition}[Classical truncation is closed]
\label{prop:classical-truncation-closed-immersion}
For every spectral scheme $X$, the canonical morphism
\[
    X_{\mathrm{cl}}
    \longrightarrow
    X
\]
is a closed immersion.
\end{proposition}

\begin{proof}
On an affine open $\Spec(A)$, the corresponding map on degree-zero rings
is the identity
\[
    \pi_0(A)
    \longrightarrow
    \pi_0(H\pi_0(A))
    \simeq
    \pi_0(A),
\]
which is surjective.
\end{proof}

\begin{definition}[Nil-immersion]
\label{def:nil-immersion-spectral}
A closed immersion $i:Z\to X$ is a \textbf{nil-immersion} if it induces an
equivalence on reduced classical schemes:
\[
    (Z_{\mathrm{cl}})_{\mathrm{red}}
    \xrightarrow{\simeq}
    (X_{\mathrm{cl}})_{\mathrm{red}}.
\]
\end{definition}

\begin{corollary}[Classical truncation is a nil-immersion]
\label{cor:classical-truncation-nil-immersion}
The canonical closed immersion
\[
    X_{\mathrm{cl}}
    \longrightarrow
    X
\]
is a nil-immersion.
\end{corollary}

\begin{proof}
It is the identity on classical truncations and therefore on reduced
classical schemes.
\end{proof}

\subsection{Open complements}

\begin{lemma}[Classical detection of the empty spectral scheme]
\label{lem:empty-spectral-scheme-classical-detection}
Let $X$ be a spectral scheme. Then
\[
    X
    \simeq
    \varnothing
\]
if and only if
\[
    X_{\mathrm{cl}}
    \simeq
    \varnothing.
\]
\end{lemma}

\begin{proof}
If $X$ is empty, then its classical truncation is empty by functoriality
of classical truncation.

Conversely, suppose
\[
    X_{\mathrm{cl}}
    \simeq
    \varnothing.
\]
Choose an affine spectral Zariski atlas
\[
    \coprod_{i\in I}
    U_i
    \longrightarrow
    X,
    \qquad
    U_i=\Spec(A_i).
\]
For every $i$, the classical truncation
\[
    (U_i)_{\mathrm{cl}}
    =
    \Spec\pi_0(A_i)
\]
is an open subscheme of $X_{\mathrm{cl}}$ and is therefore empty. Hence
\[
    \pi_0(A_i)
\]
is the zero ring.

For every integer $n$, the group $\pi_n(A_i)$ is naturally a module over
the ordinary ring $\pi_0(A_i)$. Since $\pi_0(A_i)$ is the zero ring, one
has
\[
    \pi_n(A_i)=0
\]
for every $n$. Homotopy groups detect equivalences of spectra, so
\[
    A_i
    \simeq
    0.
\]
Consequently,
\[
    U_i
    =
    \Spec(A_i)
    \simeq
    \varnothing
\]
for every $i$.

The atlas morphism is therefore an effective epimorphism
\[
    \varnothing
    \longrightarrow
    X.
\]
An effective epimorphism from the initial object can have only the
initial object as its target. Thus
\[
    X
    \simeq
    \varnothing.
\]
\end{proof}

\begin{theorem}[Complement of a closed immersion]
\label{thm:complement-closed-immersion}
Let
\[
    i
    \colon
    Z
    \hookrightarrow
    X
\]
be a closed immersion of spectral schemes. There exists an open immersion
\[
    j
    \colon
    X\setminus Z
    \hookrightarrow
    X
\]
characterized by the following universal property: a morphism
$T\to X$ factors through $X\setminus Z$ if and only if
\[
    T\times_X Z
    \simeq
    \varnothing.
\]
Its classical truncation is the ordinary open complement
\[
    (X\setminus Z)_{\mathrm{cl}}
    \simeq
    X_{\mathrm{cl}}\setminus Z_{\mathrm{cl}}.
\]
\end{theorem}

\begin{proof}
Let
\[
    U_{\mathrm{cl}}
    :=
    X_{\mathrm{cl}}
    \setminus
    Z_{\mathrm{cl}}.
\]
By
\Ref{prop:opens-classical-truncation},
this ordinary open subscheme determines a unique open spectral subscheme
\[
    j
    \colon
    U
    \hookrightarrow
    X.
\]
Define
\[
    X\setminus Z
    :=
    U.
\]
By construction,
\[
    (X\setminus Z)_{\mathrm{cl}}
    \simeq
    X_{\mathrm{cl}}
    \setminus
    Z_{\mathrm{cl}}.
\]

Let
\[
    f
    \colon
    T
    \longrightarrow
    X
\]
be a morphism of spectral schemes.

Suppose first that \(f\) factors through \(U\). Then
\[
    T\times_XZ
    \simeq
    T\times_U(U\times_XZ).
\]
The classical truncation of \(U\times_XZ\) is
\[
\begin{aligned}
    (U\times_XZ)_{\mathrm{cl}}
    &\simeq
    U_{\mathrm{cl}}
    \times_{X_{\mathrm{cl}}}
    Z_{\mathrm{cl}}
    \\
    &=
    \varnothing
\end{aligned}
\]
by
\Ref{prop:global-classical-truncation-fibre-products}.
A spectral scheme is empty if and only if its classical truncation is
empty. Hence
\[
    U\times_XZ
    \simeq
    \varnothing,
\]
and therefore
\[
    T\times_XZ
    \simeq
    \varnothing.
\]

Conversely, suppose
\[
    T\times_XZ
    \simeq
    \varnothing.
\]
Taking classical truncations and using
\Ref{prop:global-classical-truncation-fibre-products}
gives
\[
    T_{\mathrm{cl}}
    \times_{X_{\mathrm{cl}}}
    Z_{\mathrm{cl}}
    =
    \varnothing.
\]
By the ordinary universal property of the open complement, the morphism
\[
    f_{\mathrm{cl}}
    \colon
    T_{\mathrm{cl}}
    \longrightarrow
    X_{\mathrm{cl}}
\]
factors through
\[
    U_{\mathrm{cl}}
    =
    X_{\mathrm{cl}}\setminus Z_{\mathrm{cl}}.
\]

Form the pullback open immersion
\[
    j_T
    \colon
    U_T
    :=
    T\times_XU
    \longrightarrow
    T.
\]
Its classical truncation is
\[
\begin{aligned}
    (U_T)_{\mathrm{cl}}
    &\simeq
    T_{\mathrm{cl}}
    \times_{X_{\mathrm{cl}}}
    U_{\mathrm{cl}}
    \\
    &\simeq
    T_{\mathrm{cl}},
\end{aligned}
\]
where the last equivalence uses the factorization of \(f_{\mathrm{cl}}\)
through \(U_{\mathrm{cl}}\).

By
\Ref{prop:opens-classical-truncation},
open spectral subschemes of \(T\) are determined by their classical open
subschemes. Since \(U_T\to T\) has classical truncation equal to all of
\(T_{\mathrm{cl}}\), it follows that
\[
    U_T
    \xrightarrow{\simeq}
    T.
\]
The inverse equivalence gives a section
\[
    T
    \longrightarrow
    U_T
    =
    T\times_XU,
\]
and composition with the projection to \(U\) yields a factorization
\[
    T
    \longrightarrow
    U
    \longrightarrow
    X
\]
of \(f\).

The factorization is unique because
\[
    U
    \longrightarrow
    X
\]
is an open immersion and therefore a monomorphism. This proves the
universal property.
\end{proof}

\begin{proposition}[Base change of complements]
\label{prop:complement-base-change}
Let $Z\hookrightarrow X$ be closed and let $Y\to X$ be any morphism. Then
there is a canonical equivalence
\[
    Y\times_X(X\setminus Z)
    \simeq
    Y\setminus(Y\times_XZ).
\]
\end{proposition}

\begin{proof}
Both sides are open subschemes of $Y$. Their classical truncations are the
same ordinary open complement, and open spectral subschemes are determined
by their classical opens by \Ref{prop:opens-classical-truncation}.
\end{proof}

\begin{corollary}[Complement depends on the reduced closed part]
\label{cor:complement-reduced-closed-part}
If $Z\hookrightarrow X$ and $Z'\hookrightarrow X$ have the same reduced
classical closed subscheme, then
\[
    X\setminus Z
    \simeq
    X\setminus Z'.
\]
\end{corollary}

\begin{proof}
Their ordinary open complements in $X_{\mathrm{cl}}$ coincide. Apply
\Ref{prop:opens-classical-truncation}.
\end{proof}

\subsection{Pushouts along closed immersions}

\begin{theorem}[Closed pushouts]
\label{thm:closed-pushouts-spectral-schemes}
Suppose
\[
    Z
    \hookrightarrow
    X
    \qquad\text{and}\qquad
    Z
    \hookrightarrow
    Y
\]
are closed immersions of spectral schemes. Then the pushout
\[
    X\coprod_ZY
\]
exists in spectral schemes. The canonical maps
\[
    X
    \longrightarrow
    X\coprod_ZY
    \qquad\text{and}\qquad
    Y
    \longrightarrow
    X\coprod_ZY
\]
are closed immersions. The construction is compatible with arbitrary
base change.
\end{theorem}

\begin{proof}
We first record the affine calculation. Write
\[
    X=\Spec(A),
    \qquad
    Y=\Spec(B),
    \qquad
    Z=\Spec(C),
\]
so that the two closed immersions are represented by morphisms
\[
    A
    \longrightarrow
    C
    \longleftarrow
    B
\]
which are surjective on degree-zero homotopy. Put
\[
    D
    :=
    A\times_C B
\]
in $\CAlg(\Sp)$.

The forgetful functor
\[
    \CAlg(\Sp)
    \longrightarrow
    \Sp
\]
creates limits. Hence the underlying spectrum of $D$ fits into a fibre
sequence
\[
    D
    \longrightarrow
    A\oplus B
    \longrightarrow
    C,
\]
where the second morphism is the difference of the two structural
morphisms.

Since $A$, $B$, and $C$ are connective, the only possible negative
homotopy group of $D$ is $\pi_{-1}(D)$. The associated exact sequence
gives
\[
    \pi_{-1}(D)
    \simeq
    \operatorname{coker}
    \left(
        \pi_0(A)\oplus\pi_0(B)
        \longrightarrow
        \pi_0(C)
    \right).
\]
The displayed morphism is surjective because either structural
homomorphism
\[
    \pi_0(A)
    \longrightarrow
    \pi_0(C)
    \qquad\text{or}\qquad
    \pi_0(B)
    \longrightarrow
    \pi_0(C)
\]
is surjective. Therefore
\[
    \pi_{-1}(D)=0,
\]
and $D$ is connective.

The same exact sequence identifies
\[
    \pi_0(D)
    \simeq
    \pi_0(A)
    \times_{\pi_0(C)}
    \pi_0(B).
\]
The projections
\[
    \pi_0(D)
    \longrightarrow
    \pi_0(A)
    \qquad\text{and}\qquad
    \pi_0(D)
    \longrightarrow
    \pi_0(B)
\]
are surjective. Consequently,
\[
    \Spec(A)
    \longrightarrow
    \Spec(D)
    \qquad\text{and}\qquad
    \Spec(B)
    \longrightarrow
    \Spec(D)
\]
are closed immersions.

The universal property of the pullback algebra $D$ gives the universal
property of the affine pushout:
\[
    \Spec(D)
    \simeq
    \Spec(A)
    \coprod_{\Spec(C)}
    \Spec(B).
\]

We also record affine base-change compatibility. Let
\[
    D
    \longrightarrow
    D'
\]
be any morphism of connective commutative ring spectra. Tensoring the
fibre sequence
\[
    D
    \longrightarrow
    A\oplus B
    \longrightarrow
    C
\]
with $D'$ over $D$ gives a fibre sequence of $D'$-modules
\[
    D'
    \longrightarrow
    \bigl(
        A\otimes_DD'
    \bigr)
    \oplus
    \bigl(
        B\otimes_DD'
    \bigr)
    \longrightarrow
    C\otimes_DD',
\]
because extension of scalars is exact.

The compatible structural morphisms determine a canonical morphism
\[
    D'
    \longrightarrow
    \bigl(
        A\otimes_DD'
    \bigr)
    \times_{C\otimes_DD'}
    \bigl(
        B\otimes_DD'
    \bigr).
\]
The forgetful functor from commutative ring spectra to spectra creates
limits, so the underlying spectrum of the target is the fibre of
\[
    \bigl(
        A\otimes_DD'
    \bigr)
    \oplus
    \bigl(
        B\otimes_DD'
    \bigr)
    \longrightarrow
    C\otimes_DD'.
\]
The preceding fibre sequence therefore proves that the canonical
morphism is an equivalence:
\[
    D'
    \simeq
    \bigl(
        A\otimes_DD'
    \bigr)
    \times_{C\otimes_DD'}
    \bigl(
        B\otimes_DD'
    \bigr).
\]
Passing to affine spectra proves arbitrary base-change compatibility in
the affine case.

For general spectral schemes, the affine-local notion of closed
immersion used in this chapter agrees with the spectral-scheme notion of
closed immersion in the closed-gluing theorem; see
\cite[Theorem~4.4 and Remark~4.5]{LurieDAGIX}.
The spectral-scheme form of the closed-gluing theorem gives a pushout
diagram
\[
\begin{tikzcd}[column sep=large, row sep=large]
    Z
        \arrow[r, hook]
        \arrow[d, hook]
    &
    X
        \arrow[d, hook]
    \\
    Y
        \arrow[r, hook]
    &
    P
\end{tikzcd}
\]
in spectral schemes, and asserts that the two canonical morphisms into
$P$ are closed immersions; see
\cite[Theorem~6.1 and Remark~6.2]{LurieDAGIX}.
Thus
\[
    P
    \simeq
    X\coprod_ZY.
\]

Finally, the base-change theorem for closed gluing asserts that, for
every morphism
\[
    T
    \longrightarrow
    P,
\]
the induced square
\[
\begin{tikzcd}[column sep=large, row sep=large]
    T\times_PZ
        \arrow[r, hook]
        \arrow[d, hook]
    &
    T\times_PX
        \arrow[d, hook]
    \\
    T\times_PY
        \arrow[r, hook]
    &
    T
\end{tikzcd}
\]
is again a pushout square. The same proof applies to spectral schemes;
see \cite[Proposition~9.2]{LurieDAGIX}. Hence the construction is
compatible with arbitrary base change.
\end{proof}

\begin{proposition}[Pushouts of nil-immersions]
\label{prop:pushout-nil-immersions}
In the situation of \Ref{thm:closed-pushouts-spectral-schemes}, if one of
the two closed immersions is a nil-immersion, then the opposite canonical
map into the pushout is a nil-immersion.
\end{proposition}
\begin{proof}
The assertion is local on the pushout, so it is enough to prove the
affine statement. Write
\[
    X=\Spec(A),
    \qquad
    Y=\Spec(B),
    \qquad
    Z=\Spec(C),
\]
and let
\[
    D
    :=
    A\times_CB.
\]
By the affine calculation in the proof of
\Ref{thm:closed-pushouts-spectral-schemes},
the pushout is
\[
    X\coprod_ZY
    \simeq
    \Spec(D),
\]
and
\[
    \pi_0(D)
    \simeq
    \pi_0(A)
    \times_{\pi_0(C)}
    \pi_0(B).
\]

Suppose first that
\[
    Z
    \longrightarrow
    X
\]
is a nil-immersion. Since it is a closed immersion, the map
\[
    \pi_0(A)
    \twoheadrightarrow
    \pi_0(C)
\]
is surjective. The nil-immersion condition says that its kernel
\[
    I
    :=
    \ker
    \bigl(
        \pi_0(A)
        \longrightarrow
        \pi_0(C)
    \bigr)
\]
is a nil ideal.

The projection
\[
    \pi_0(D)
    =
    \pi_0(A)
    \times_{\pi_0(C)}
    \pi_0(B)
    \longrightarrow
    \pi_0(B)
\]
is surjective, and its kernel is canonically identified with \(I\):
\[
    \ker
    \bigl(
        \pi_0(D)
        \longrightarrow
        \pi_0(B)
    \bigr)
    \simeq
    I.
\]
This kernel is therefore a nil ideal. It follows that the closed
immersion
\[
    Y
    =
    \Spec(B)
    \longrightarrow
    \Spec(D)
    =
    X\coprod_ZY
\]
induces an equivalence on reduced classical schemes. Hence it is a
nil-immersion.

The argument with \(A\) and \(B\) interchanged proves that if
\[
    Z
    \longrightarrow
    Y
\]
is a nil-immersion, then
\[
    X
    \longrightarrow
    X\coprod_ZY
\]
is a nil-immersion. The affine conclusions glue because nil-immersions
are detected on classical reductions and are local on the target.
\end{proof}

\section{Vector bundles and brave new affine spaces}
\label{sec:vector-bundles-brave-new-affine-spaces}

\subsection{Locally free quasi-coherent modules}

\begin{definition}[Locally free quasi-coherent module]
\label{def:locally-free-qcoh}
Let $X$ be a spectral scheme. A quasi-coherent module
$\mathcal E\in\QCoh(X)$ is \textbf{locally free of rank $n$} if there is a
Zariski open cover $\{U_i\to X\}$ such that
\[
    \mathcal E|_{U_i}
    \simeq
    \mathcal O_{U_i}^{\oplus n}
\]
for every $i$.
\end{definition}

A \textbf{vector bundle of rank $n$} on $X$ means a locally free
quasi-coherent module of rank $n$.

\begin{proposition}[Locally free modules are finite projective]
\label{prop:locally-free-finite-projective}
Let
\[
    \mathcal E
    \in
    \QCoh(X)
\]
be locally free of rank \(n<\infty\). Then \(\mathcal E\) is dualizable,
connective, and flat. On every affine open it corresponds to a finite
projective module.
\end{proposition}

\begin{proof}
Choose a Zariski open cover
\[
    \{U_i\to X\}_{i\in I}
\]
such that
\[
    \mathcal E|_{U_i}
    \simeq
    \mathcal O_{U_i}^{\oplus n}
\]
for every \(i\).

Each restriction \(\mathcal E|_{U_i}\) is dualizable, with dual
\[
    (\mathcal E|_{U_i})^\vee
    \simeq
    \mathcal O_{U_i}^{\oplus n}.
\]
The local duals, evaluation maps, and coevaluation maps are compatible
on overlaps. By
\Ref{thm:qcoh-atlas-computation},
they glue to a quasi-coherent module
\[
    \mathcal E^\vee
    \in
    \QCoh(X)
\]
and morphisms
\[
    \mathcal E^\vee
    \otimes_{\mathcal O_X}
    \mathcal E
    \longrightarrow
    \mathcal O_X
\]
and
\[
    \mathcal O_X
    \longrightarrow
    \mathcal E
    \otimes_{\mathcal O_X}
    \mathcal E^\vee.
\]
The triangle identities hold after restriction to the chosen open cover
and therefore hold globally. Thus \(\mathcal E\) is dualizable.

Connectivity is detected on affine opens by the canonical \(t\)-structure
of
\Ref{thm:qcoh-spectral-scheme-t-structure}.
Since every restriction
\[
    \mathcal E|_{U_i}
    \simeq
    \mathcal O_{U_i}^{\oplus n}
\]
is connective, \(\mathcal E\) is connective.

We next prove flatness. Let
\[
    U
    =
    \Spec(A)
    \subseteq
    X
\]
be an affine open, and let
\[
    M
    \in
    \Mod_A
\]
correspond to \(\mathcal E|_U\). Pulling the chosen trivializing cover
back to \(U\) and refining by affine opens gives an affine spectral
Zariski covering family
\[
    \{\Spec(A_j)\to\Spec(A)\}_{j\in J}
\]
such that
\[
    M\otimes_AA_j
    \simeq
    A_j^{\oplus n}
\]
for every \(j\). Each finite free \(A_j\)-module is flat. By
\Ref{prop:flat-module-zariski-local},
the \(A\)-module \(M\) is flat. Since this holds on every affine open
\(U\subseteq X\), the quasi-coherent module \(\mathcal E\) is flat.

Finally, let
\[
    U
    =
    \Spec(A)
    \subseteq
    X
\]
be any affine open, and let \(M\) denote the corresponding \(A\)-module.
The restriction \(\mathcal E|_U\) is dualizable, so \(M\) is perfect by
\Ref{prop:perfect-equals-dualizable}.
The preceding arguments show that \(M\) is connective and flat.
Therefore
\Ref{prop:finite-projective-characterization}
implies that \(M\) is finite projective over \(A\).
\end{proof}

\begin{proposition}[Base change of locally free modules]
\label{prop:base-change-locally-free}
If $f:Y\to X$ is any morphism and $\mathcal E$ is locally free of rank
$n$, then $f^*\mathcal E$ is locally free of rank $n$.
\end{proposition}

\begin{proof}
Pull back a trivializing Zariski cover of $X$. Open covers and finite free
modules are preserved under base change.
\end{proof}

\subsection{The relative free commutative algebra}

\begin{proposition}[Relative free commutative algebra]
\label{prop:relative-free-commutative-algebra}
For every spectral scheme $X$, the forgetful functor
\[
    U_X
    \colon
    \operatorname{QCAlg}(X)
    \longrightarrow
    \QCoh(X)
\]
admits a left adjoint
\[
    \operatorname{Sym}_{\mathcal O_X}
    \colon
    \QCoh(X)
    \longrightarrow
    \operatorname{QCAlg}(X).
\]
The construction is compatible with pullback: for every $f:Y\to X$ there
is a canonical equivalence
\[
    f^*
    \operatorname{Sym}_{\mathcal O_X}(\mathcal E)
    \simeq
    \operatorname{Sym}_{\mathcal O_Y}(f^*\mathcal E).
\]
\end{proposition}

\begin{proof}
The category $\QCoh(X)$ is presentably symmetric monoidal, so the
operadic free-algebra theorem supplies the adjunction. Since pullback is
colimit-preserving and strong symmetric monoidal, it commutes with the
operadic free commutative algebra construction.
\end{proof}

\begin{proposition}[Connectivity of the relative free algebra]
\label{prop:relative-free-algebra-connective}
If
\[
    \mathcal E
    \in
    \QCoh(X)_{\geq0},
\]
then
\[
    \operatorname{Sym}_{\mathcal O_X}(\mathcal E)
    \in
    \operatorname{QCAlg}(X)^{\mathrm{cn}}.
\]
\end{proposition}

\begin{proof}
The assertion is affine-local. Let $U=\Spec(A)\subseteq X$, and let
$M\in\Mod_{A,\geq0}$ correspond to $\mathcal E|_U$. The underlying
$A$-module of the free commutative $A$-algebra is
\[
    \bigoplus_{k\geq0}
    \left(
        M^{\otimes_Ak}
    \right)_{h\Sigma_k}.
\]
Every tensor power $M^{\otimes_Ak}$ is connective. Homotopy orbits and
small coproducts preserve connective objects, so the displayed
$A$-module is connective. Connectivity of a quasi-coherent algebra is
detected on affine opens, and the result follows.
\end{proof}

\subsection{Total spaces and affine spaces}

\begin{definition}[Linear spectral space]
\label{def:linear-spectral-space}
For a connective quasi-coherent module
$\mathcal E\in\QCoh(X)_{\geq0}$, define its \textbf{linear spectral
space} by
\[
    \mathbb V_X(\mathcal E)
    :=
    \Spec_X
    \operatorname{Sym}_{\mathcal O_X}(\mathcal E).
\]
\end{definition}

\begin{proposition}[Functor of points of a linear spectral space]
\label{prop:linear-space-functor-points}
For every morphism $f:T\to X$, there is a natural equivalence
\[
    \Map_X
    \bigl(
        T,
        \mathbb V_X(\mathcal E)
    \bigr)
    \simeq
    \Map_{\QCoh(T)}
    \bigl(
        f^*\mathcal E,
        \mathcal O_T
    \bigr).
\]
\end{proposition}

\begin{proof}
Apply the relative spectrum universal property of
\Ref{thm:relative-spectrum-construction} and the free-algebra adjunction of
\Ref{prop:relative-free-commutative-algebra}.
\end{proof}

\begin{remark}[Dual convention]
\label{rem:vector-bundle-dual-convention}
With the convention of \Ref{def:linear-spectral-space}, points correspond
to linear maps from $\mathcal E$ to the structure sheaf. If one wishes the
points to correspond to sections of a finite locally free module
$\mathcal E$, one uses
\[
    \mathbf V_X(\mathcal E)
    :=
    \mathbb V_X(\mathcal E^\vee).
\]
The brave new affine spaces below use the self-dual free modules, so the
two conventions agree canonically there.
\end{remark}

\begin{definition}[Brave new affine space]
\label{def:brave-new-affine-space}
For an integer $n\geq0$, define
\[
    \mathbb A_X^n
    :=
    \mathbb V_X
    \bigl(
        \mathcal O_X^{\oplus n}
    \bigr)
    =
    \Spec_X
    \operatorname{Sym}_{\mathcal O_X}
    \bigl(
        \mathcal O_X^{\oplus n}
    \bigr).
\]
The case $n=1$ is the \textbf{brave new affine line}
\[
    \mathbb A_X^1
    =
    \Spec_X
    \operatorname{Sym}_{\mathcal O_X}(\mathcal O_X).
\]
\end{definition}

\begin{proposition}[Base change of affine space]
\label{prop:affine-space-base-change}
For every morphism $f:Y\to X$, there is a canonical equivalence
\[
    \mathbb A_X^n
    \times_X
    Y
    \simeq
    \mathbb A_Y^n.
\]
\end{proposition}

\begin{proof}
Combine the pullback compatibility of the free algebra from
\Ref{prop:relative-free-commutative-algebra} with relative spectrum base
change from \Ref{prop:relative-spectrum-base-change}.
\end{proof}

\begin{proposition}[Classical truncation of affine space]
\label{prop:classical-truncation-affine-space}
There is a canonical equivalence of ordinary schemes
\[
    (\mathbb A_X^n)_{\mathrm{cl}}
    \simeq
    \mathbb A_{X_{\mathrm{cl}}}^n.
\]
\end{proposition}

\begin{proof}
The assertion is local on $X$. For $X=\Spec(A)$, the homotopy-degree-zero
part of the free commutative $A$-algebra on $A^{\oplus n}$ is the ordinary
polynomial algebra
\[
    \pi_0(A)[t_1,\ldots,t_n].
\]
Taking affine spectra gives the result.
\end{proof}

\begin{construction}[Zero and unit sections]
\label{con:zero-unit-sections-affine-line}
The zero endomorphism and identity endomorphism of $\mathcal O_X$ determine
maps of quasi-coherent modules
\[
    \mathcal O_X
    \rightrightarrows
    \mathcal O_X.
\]
By the free-algebra and relative-spectrum adjunctions, they induce sections
\[
    0,1
    \colon
    X
    \rightrightarrows
    \mathbb A_X^1.
\]
\end{construction}

\begin{proposition}[The zero section is closed]
\label{prop:zero-section-affine-line-closed}
The zero section
\[
    0
    \colon
    X
    \longrightarrow
    \mathbb A_X^1
\]
is a closed immersion.
\end{proposition}

\begin{proof}
The assertion is local on $X$. Let
\[
    X=\Spec(A).
\]
Then
\[
    \mathbb A_X^1
    \simeq
    \Spec
    \operatorname{Sym}_A(A),
\]
and the zero section is represented by the augmentation
\[
    \epsilon_0
    \colon
    \operatorname{Sym}_A(A)
    \longrightarrow
    A
\]
which sends the universal generator to zero.

On degree-zero homotopy groups, this is the ordinary polynomial
augmentation
\[
    \pi_0(A)[t]
    \longrightarrow
    \pi_0(A),
    \qquad
    t
    \longmapsto
    0.
\]
It is surjective. Therefore the zero section is a closed immersion by
\Ref{def:affine-closed-immersion-spectral}.
\end{proof}

\begin{definition}[Spectral multiplicative group]
\label{def:spectral-multiplicative-group}
The \textbf{spectral multiplicative group} over $X$ is the open complement
\[
    \mathbb G_{m,X}
    :=
    \mathbb A_X^1
    \setminus
    0(X).
\]
\end{definition}

\begin{proposition}[Group structure on the spectral multiplicative group]
\label{prop:spectral-multiplicative-group-structure}
The spectral scheme
\[
    \mathbb G_{m,X}
\]
is canonically a commutative group object in spectral schemes over $X$.

For every morphism $T\to X$, there is a natural equivalence
\[
    \Map_X
    \bigl(
        T,
        \mathbb G_{m,X}
    \bigr)
    \simeq
    \operatorname{Aut}_{\QCoh(T)}
    (\mathcal O_T),
\]
where the right side denotes the space of invertible
$\mathcal O_T$-linear endomorphisms of $\mathcal O_T$.
\end{proposition}

\begin{proof}
The assertion is local on $X$. Let $X=\Spec(A)$ and write
\[
    A\{t\}
    :=
    \operatorname{Sym}_A(A).
\]
The complement of the zero section is the principal open where the
universal element $t$ is invertible:
\[
    \mathbb G_{m,X}
    \simeq
    \Spec
    \bigl(
        A\{t\}[t^{-1}]
    \bigr).
\]

The coordinate algebra carries the commutative Hopf algebra maps
\[
    \Delta
    \colon
    A\{t\}[t^{-1}]
    \longrightarrow
    A\{t_1,t_2\}
    [(t_1t_2)^{-1}],
    \qquad
    t\longmapsto t_1t_2,
\]
\[
    \epsilon
    \colon
    A\{t\}[t^{-1}]
    \longrightarrow
    A,
    \qquad
    t\longmapsto1,
\]
and
\[
    \iota
    \colon
    A\{t\}[t^{-1}]
    \longrightarrow
    A\{t\}[t^{-1}],
    \qquad
    t\longmapsto t^{-1}.
\]
The universal property of localization gives these maps uniquely, and
the commutative group identities follow from the corresponding identities
satisfied by the universal invertible element. Passing contravariantly to
relative spectra defines multiplication, unit, and inversion on
$\mathbb G_{m,X}$.

For $T\to X$, a map
\[
    T
    \longrightarrow
    \mathbb A_X^1
\]
corresponds to an $\mathcal O_T$-linear endomorphism of $\mathcal O_T$.
It factors through the principal open $\mathbb G_{m,X}$ exactly when the
corresponding element is invertible. This gives the displayed
functor-of-points equivalence.

The construction is compatible with affine restriction and therefore
glues over an arbitrary spectral scheme $X$.
\end{proof}

\begin{proposition}[Base change of the multiplicative group]
\label{prop:multiplicative-group-base-change}
For every morphism $Y\to X$, there is a canonical equivalence
\[
    \mathbb G_{m,X}
    \times_X
    Y
    \simeq
    \mathbb G_{m,Y}.
\]
\end{proposition}

\begin{proof}
Use affine-space base change and base change of open complements from
\Ref{prop:complement-base-change}.
\end{proof}

\begin{remark}[The brave new affine line is genuinely spectral]
\label{rem:affine-line-genuinely-spectral}
Even when $X$ is classical, the free commutative $\EE_\infty$-algebra
\[
    \operatorname{Sym}_{\mathcal O_X}(\mathcal O_X)
\]
need not be discrete. Thus $\mathbb A_X^1$ is not obtained merely by
viewing the ordinary affine line as a spectral scheme. Its classical
truncation is ordinary affine space by
\Ref{prop:classical-truncation-affine-space}, but its higher spectral
structure is part of the intended motivic interval. See
\cite{KhanBraveNewMotivicI}.
\end{remark}



\section{The spectral-geometric output}
\label{sec:spectral-geometric-output}

\begin{theorem}[Spectral algebraic geometry package]
\label{thm:spectral-algebraic-geometry-package}
The constructions of this chapter canonically provide:
\begin{enumerate}[label=(\roman*)]
    \item the $\infty$-category $\operatorname{Sch}^{\mathrm{sp}}$ of
    spectral schemes and its full subcategory of affine schemes
    $\Aff_{\Sp}$;

    \item classical truncation
    \[
        (-)_{\mathrm{cl}}
        \colon
        \operatorname{Sch}^{\mathrm{sp}}
        \longrightarrow
        \operatorname{Sch},
    \]
    compatible with fibre products and detecting open subsets,
    affineness, quasi-compactness, and quasi-separatedness;

    \item the symmetric monoidal quasi-coherent coefficient system
    $X\mapsto\QCoh(X)$, its canonical $t$-structure on spectral schemes,
    and the relative spectrum functor
    \[
        \Spec_X
        \colon
        \operatorname{QCAlg}(X)^{\mathrm{cn},\op}
        \longrightarrow
        \operatorname{Sch}^{\mathrm{sp}}_{/X};
    \]

    \item flat morphisms, open and closed immersions, open complements,
    pushouts along closed immersions, affine morphisms, vector bundles,
    relative linear spaces, affine spaces, and the spectral
    multiplicative group, all with the base-change properties established
    above.
\end{enumerate}
These constructions form the geometric base for the infinitesimal theory
of \Ref{chap:spectral-deformation-theory}.
\end{theorem}

\begin{proof}
The assertions collect the constructions and permanence theorems of
\Ref{sec:affine-spectral-schemes} through
\Ref{sec:vector-bundles-brave-new-affine-spaces}.
\end{proof}
